%%%%%%%%%%%%%%%%%%%%%%%%%%%%%%%%%%%%%%%%%%%%%%%%%%%%%%%%%%%%%%%%%%%%%%%%%%%%%%
%  Marketing Management (IM345AT) -- Volume 5 of 5
%  UNIT V : Conversion Optimization and Digital Analytics
%%%%%%%%%%%%%%%%%%%%%%%%%%%%%%%%%%%%%%%%%%%%%%%%%%%%%%%%%%%%%%%%%%%%%%%%%%%%%%

\documentclass{MMTextbook}

\unitnumber{V}
\unittitle{Conversion Optimization and Digital Analytics}
\unitsubtitle{Turning attention into action, and turning measurement into
decisions}
\unithours{07 Hours}
\volumeno{5}

\makeindex[intoc=true,columns=2,title=Index]

\hypersetup{
  pdftitle={Marketing Management IM345AT -- Unit V: Conversion Optimization and
             Digital Analytics},
  pdfsubject={Semester IV, Department of Industrial Engineering and Management,
              RV College of Engineering},
  pdfkeywords={conversion optimization, AIDAS, landing page, social proof,
               digital analytics, IM345AT, marketing management}
}

\begin{document}

%%%%%%%%%%%%%%%%%%%%%%%%%%%%%%%%%%%%%%%%%%%%%%%%%%%%%%%%%%%%%%%%%%%%%%%%%%%%%%
\frontmatter
\pagestyle{mmfront}

\halftitlepage
\maintitlepage

\colophonpage{%
\textbf{\coursetitle\ ---\ The Study Text}\\[2pt]
Volume 5 of 5 \textbullet\ Unit V: Conversion Optimization and Digital
Analytics\\[10pt]
%
Prepared for the course \coursetitle\ (\coursecode), Semester IV, of the
\coursedept, \courseinstitute.\\[10pt]
%
\textbf{Sources.} The syllabus, course outcomes, assessment rubrics and
reference-book list reproduced in the front matter are transcribed from the
official course document. The previous-year questions in Appendix~B are
transcribed from RVCE Continuous Internal Evaluation and Semester End
Examination papers for this course, and where an official scheme of valuation
was released the model answers follow it. The explanatory text is compiled from
classroom lecture notes and, where those were incomplete, from the prescribed
reference texts.\\[10pt]
%
\textbf{Typography.} Set in Nimbus Roman with URW Gothic display headings.
Composed in \LaTeX{} using the \texttt{MMTextbook} class written for this
series. The visual design is derived from \emph{The Legrand Orange Book}
template by Mathias Legrand, with modifications by Vel, distributed through
LaTeXTemplates.com under a Creative Commons BY-NC-SA 4.0 licence; the class
re-implements that design language on standard CTAN packages. All figures are
drawn in \TikZ{} and \texttt{pgfplots}.\\[10pt]
%
\textbf{Status.} Unofficial student-compiled study material. Not issued,
endorsed or verified by the Department or by RV College of Engineering.\\[10pt]
%
2026 edition.}

% ---------------------------------------------------------------- preface
\frontchapter{Preface to Volume 5}

Unit~V closes the course by answering the two questions everything before it has
been building towards: \emph{why did that visitor not buy?}, and \emph{how does
the person holding the data get the organisation to act on it?}

The first half is conversion optimization, and it is organised around one model.
AIDAS --- Attention, Interest, Desire, Action, Satisfaction --- is the single most
examined topic in this unit; it has appeared in every Unit~V paper available, as
both a two-mark definition and a ten-mark explanation. It is worth learning not as
a list of five words but as a \emph{diagnostic}: each stage corresponds to a
distinct failure with a distinct symptom, so knowing the model tells you which
element of a page to fix.

The second half is digital analytics, and its centre of gravity is not
measurement but influence. The syllabus is explicit on this point: after the
evolution of the discipline and the end-to-end customer experience, it names the
analyst's \emph{influence on business} and \emph{role as a change agent}. That is
unusual for a technical topic, and it is deliberate. An organisation with
excellent tools and poor interpretation makes confident wrong decisions faster
than one with no tools at all.

Chapters~1 to~4 cover conversion optimization; Chapters~5 to~7 cover digital
analytics. The chapter sequence follows the syllabus phrase order exactly, and
Appendix~A records the mapping.

Two notes specific to this volume.

First, Chapter~3 --- what visitors want to see --- deliberately repeats material
in Volume~4. Both unit descriptors name it: Unit~IV under ``User-centered Design
and Website Optimization'' and Unit~V under ``what visitors want to see on the
website''. Volume~4 treats it as a design brief; this volume treats it as a
conversion diagnosis.

Second, this volume needs no appendix of off-syllabus material. Every topic in
the source notes maps to a descriptor phrase.

\vspace{12pt}
\noindent\emph{The whole of the course syllabus is reproduced overleaf. The
unit covered by this volume is Unit~V.}

% ------------------------------------------------- conventions & syllabus
%%%%%%%%%%%%%%%%%%%%%%%%%%%%%%%%%%%%%%%%%%%%%%%%%%%%%%%%%%%%%%%%%%%%%%%%%%%%%%
%  mm-howtouse.tex
%  Editorial conventions for the series.  Shared, identical, across all five
%  volumes so that a reader who learns the apparatus once can use any volume.
%%%%%%%%%%%%%%%%%%%%%%%%%%%%%%%%%%%%%%%%%%%%%%%%%%%%%%%%%%%%%%%%%%%%%%%%%%%%%%

\frontchapter{How to Use This Study Text}

\section*{The shape of a volume}
\addcontentsline{toc}{section}{The shape of a volume}

Each of the five volumes in this series covers exactly one unit of the
IM345AT syllabus and is self-contained. Every volume opens with the official
syllabus for the whole course, so that the unit in hand can always be located
within it, and closes with a set of appendices carrying the concordance,
the question bank, the glossary and a quick-reference sheet.

The numbered chapters of each volume map one-to-one onto the topic phrases of
that unit's official syllabus descriptor. The order of the chapters is the
order of the syllabus. Appendix~A sets out that mapping explicitly, phrase by
phrase, so coverage can be audited rather than assumed.

\section*{The five panel types}
\addcontentsline{toc}{section}{The five panel types}

Material that is not continuous prose is set in one of five panels. Each has a
fixed meaning throughout the series.

\begin{definitionbox}[the term being defined]
A formal definition, worth reproducing close to verbatim in an answer script.
The term appears in the panel heading and again in the index.
\end{definitionbox}

\begin{keybox}[High Yield]
A result, list or distinction that carries disproportionate examination weight.
Where a topic has recurred across papers, the recurrence is recorded as a
bracketed tag. \pyq{CIE-I, 3 Apr 2025; CIE-I, 16 Apr 2026}
\end{keybox}

\begin{exambox}[Model Answer]
An answer written at the length and in the register the marking scheme
rewards. Where an official scheme of valuation was released, the panel follows
it.
\end{exambox}

\begin{examplebox}[Worked Example]
A concrete illustration --- a named company, a named persona, a worked
calculation --- rather than a general statement.
\end{examplebox}

\begin{warnbox}
A common and costly error: a distinction usually collapsed, a term usually
misused, or an answer that looks complete and is not.
\end{warnbox}

A sixth panel, in green, sets out a framework or process as a single
sequence:

\begin{frameworkbox}[Framework]
Step one \ra\ step two \ra\ step three \ra\ back to step one.
\end{frameworkbox}

\section*{How previous-year questions are recorded}
\addcontentsline{toc}{section}{How previous-year questions are recorded}

\begin{keybox}[The Bracket Convention]
A past-paper appearance is recorded \textbf{only} as a bracketed source tag
attached to the material, never as a sentence of narrative. So the text reads
\begin{quote}\sffamily\footnotesize
Push versus pull marketing. \pyq{CIE-I, 3 Apr 2025; CIE-I, 16 Apr 2026; SEE, Jun/Jul 2025 --- 2 M}
\end{quote}
and never ``this was asked in 2025 and again in 2026''. The tag carries the
paper, the date and, where the mark allocation is known, the marks. Multiple
appearances are separated by semicolons.
\end{keybox}

\noindent The abbreviations used inside tags are:

\vspace{4pt}
{\sffamily\small
\begin{tabular}{@{}P{0.14\textwidth}P{0.80\textwidth}@{}}
\toprule[1.1pt]
\rlab{CIE-I} & First Continuous Internal Evaluation test\\
\rlab{CIE-II} & Second Continuous Internal Evaluation test\\
\rlab{CIE-III} & Third Continuous Internal Evaluation test\\
\rlab{SEE} & Semester End Examination\\
\rlab{2 M} & Mark allocation of the question as printed on the paper\\
\bottomrule[1.1pt]
\end{tabular}}

\vspace{10pt}

Appendix~B of each volume reproduces every past-paper question traced to that
unit, in full, grouped by mark weight, each carrying the same bracketed tag.

\section*{Cross-references, figures and the index}
\addcontentsline{toc}{section}{Cross-references, figures and the index}

Cross-references are given as \S\,chapter.section --- so \S\,1.4 is the fourth
section of Chapter~1. Every figure in this series is drawn in \TeX{} rather
than imported as an image, so the diagrams print at full resolution at any
size and use the same typeface as the surrounding text. Figures and tables are
listed after the table of contents.

Terms set in \textbf{\textit{bold italic}} on first use are indexed. The index
at the back of each volume is a real back-of-book index built from those
terms, not a repeat of the contents list.

\section*{Status of this text}
\addcontentsline{toc}{section}{Status of this text}

\begin{warnbox}[Unofficial Document]
This is a student-compiled study text. It is not issued, endorsed or verified
by the Department of Industrial Engineering and Management or by RV College of
Engineering. The syllabus reproduced in the front matter, the assessment
rubrics and the previous-year questions are transcribed from official
documents; the explanatory prose around them is not official. Where the source
lecture notes were incomplete, gaps have been filled from the prescribed
reference texts listed in the front matter. Always cross-check against the
official syllabus and your own class notes.
\end{warnbox}

%%%%%%%%%%%%%%%%%%%%%%%%%%%%%%%%%%%%%%%%%%%%%%%%%%%%%%%%%%%%%%%%%%%%%%%%%%%%%%
%  mm-syllabus.tex
%  Verbatim reproduction of the official IM345AT course document.
%  Shared, byte-identical, across all five volumes of the series.
%%%%%%%%%%%%%%%%%%%%%%%%%%%%%%%%%%%%%%%%%%%%%%%%%%%%%%%%%%%%%%%%%%%%%%%%%%%%%%

\frontchapter{Official Syllabus: IM345AT}

\noindent
{\sffamily\small\color{slatelight}%
Reproduced from the official course document issued by the Department of
Industrial Engineering and Management, RV College of Engineering. The text of
the five unit descriptors below is the authority against which every chapter of
this series is written; nothing outside it appears in the numbered chapters of
any volume.\par}

\vspace{10pt}

\begin{center}
\sffamily\small
\begin{tabular}{@{}P{0.238\textwidth}P{0.238\textwidth}P{0.218\textwidth}P{0.218\textwidth}@{}}
\toprule[1.1pt]
\multicolumn{4}{@{}l}{\bfseries\normalsize Semester: IV \quad\textbullet\quad MARKETING MANAGEMENT}\\
\multicolumn{4}{@{}l}{\color{slatelight}Category: Professional Core Course (Theory)}\\
\midrule
\rlab{Course Code} & IM345AT & \rlab{CIE Marks} & 100 Marks\\
\rlab{Credits (L:T:P)} & 3:0:0 & \rlab{SEE Marks} & 100 Marks\\
\rlab{Total Hours} & 45 L & \rlab{SEE Duration} & 3.00 Hours\\
\bottomrule[1.1pt]
\end{tabular}
\end{center}

\vspace{6pt}

% ---------------------------------------------------------------- unit blocks
\newcommand{\syllunit}[3]{%
  \begin{tcolorbox}[mmbase, colback=white, colframe=ocre!45, boxrule=0.7pt,
      sharp corners, fontupper=\small,
      fonttitle=\sffamily\bfseries\small\color{white},
      coltitle=white, colbacktitle=#3,
      title={#1},
      attach boxed title to top left={xshift=-0.7pt, yshift=-0.7pt},
      boxed title style={sharp corners, boxrule=0pt, colback=#3},
      top=16pt]
    #2
  \end{tcolorbox}}

\syllunit{UNIT--I \hfill 07 Hrs}{%
\textbf{Introduction to Digital Marketing:} Principles of Digital Marketing;
Digital Marketing Channels; Tools to Create Buyer Persona; Competitor Research
Tools, Website Analysis Tools, etc.\par
\vspace{3pt}
\textbf{Content Marketing:} Content Marketing Concepts \& Strategies; Planning,
Creating, Distributing \& Promoting Content; Optimize Website UX \& Landing
Pages; Measure Impact; Metrics \& Performance; Using Content Research for
Opportunities, etc.}{ocre}

\syllunit{UNIT--II \hfill 08 Hrs}{%
\textbf{Social Media Marketing:} Introduction; Major Social Media Platforms for
Marketing; Developing Data-driven Audience \& Campaign Insights; Social Media
for Business; Creation \& Optimization of Social Media Campaigns, etc.\par
\vspace{3pt}
\textbf{Search Engine Optimization:} Search Engine Optimization Fundamentals;
Keywords and SEO Content Plan; SEO \& Business Objectives; Writing SEO Content;
On-site \& off-site SEO; Optimize Organic Search Ranking, etc.}{ocre}

\syllunit{UNIT--III \hfill 07 Hrs}{%
\textbf{Web Analytics \& Google Analytics:} Google Analytics Tools; Web
Analytics Tools, etc.\par
\vspace{3pt}
\textbf{E-mail Marketing:} Effective E-mail Campaigns; E-mail Plan; E-mail
Marketing Campaign Analysis; Measuring Conversions \& keeping up, etc.}{ocre}

\syllunit{UNIT--IV \hfill 07 Hrs}{%
\textbf{Web Design:} Web design, optimization of websites; Publishing a basic
website; User-centered Design and Website Optimization; Design Principles and
Website Copy; Website Metrics \& Developing Insight, etc.\par
\vspace{3pt}
\textbf{Mobile Marketing:} Difference between mobile advertising and marketing,
utilizing mobile marketing for sales promotions, online applications,
etc.}{ocre}

\syllunit{UNIT--V \hfill 07 Hrs}{%
\textbf{Conversion Optimization:} What is AIDAS and its role; website
optimization; what visitors want to see on the website; how to optimize key
element and increase the effect of landing on a particular page\par
\vspace{3pt}
\textbf{Digital Analytics:} Evolution of Digital Analytics, information about
end-to-end customer experience, analyst's influence on business, role as a
change agent, etc.}{ocre}

% ------------------------------------------------------------ course outcomes
\section*{Course Outcomes}
\addcontentsline{toc}{section}{Course Outcomes}

\noindent{\small After completing the course, the students will be able to:\par}
\vspace{4pt}

{\sffamily\small
\begin{longtable}{@{}P{0.07\textwidth}P{0.89\textwidth}@{}}
\toprule[1.1pt]
\rlab{CO1} & Differentiate the benefits drawn by updated marketing mix from
traditional marketing mix for effective marketing management there by to stay
competitive in today's global market-place.\\
\midrule
\rlab{CO2} & Develop an effective holistic marketing atmosphere to efficiently
face the challenges in dynamically changing market.\\
\midrule
\rlab{CO3} & Formulate a potential marketing plan to effectively reach the
targeted market segments, by delivering the value to targeted customers through
practicing sound marketing research.\\
\midrule
\rlab{CO4} & Create new channels to improvise marketing to achieve and maintain
competitive position in globalized market-place.\\
\bottomrule[1.1pt]
\end{longtable}}

% --------------------------------------------------------------- reference books
\section*{Prescribed Reference Books}
\addcontentsline{toc}{section}{Prescribed Reference Books}

{\small
\begin{enumerate}[leftmargin=2.1em]
\item \textit{Marketing Management}, Philip Kotler, Kevin Lane Keller, 15th
      Edition, 2016, Pearson, ISBN: 978-93-325-5718-5.
\item \textit{Digital Marketing --- Strategy, Implementation \& Practice}, Dave
      Chaffey, Fiona Ellis-Chadwick, 7th Edition, 2019, Pearson,
      ISBN: 9781292241623, 1292241624.
\item \textit{Marketing Research}, Donald S. Tull, Del I. Hawkins, 6th Edition,
      Prentice Hall India, ISBN: 8120309618.
\item \textit{Marketing Management --- A South Asian Perspective}, Philip
      Kotler, Kevin Lane Keller, Abraham Koshy, Mithileshwar Jha, 14th Edition,
      2013, Pearson, ISBN: 978-81-317-6716-0.
\item \textit{Marketing Research}, David A. Aaker, V. Kumar, George S. Day, 9th
      Edition, 2008, John Wiley \& Sons, ISBN: 978-265-1791-6.
\end{enumerate}}

% ------------------------------------------------------------------- rubrics
\frontchapter{Assessment Scheme}

\section*{Rubric for the Continuous Internal Evaluation (Theory)}
\addcontentsline{toc}{section}{Rubric for the Continuous Internal Evaluation}

{\sffamily\small
\begin{longtable}{@{}P{0.045\textwidth}P{0.80\textwidth}>{\raggedleft\arraybackslash}P{0.07\textwidth}@{}}
\toprule[1.1pt]
\rowcolor{ocre}\thd{\#} & \thd{Component} & \thd{Marks}\\
\midrule
\rlab{1.} & \textbf{Quizzes:} Quizzes will be conducted in online/offline mode.
TWO QUIZZES will be conducted \& each quiz will be evaluated for 10 marks. The
sum of two quizzes will be the final quiz marks. & \bfseries 20\\
\midrule
\rowcolor{tablealt}
\rlab{2.} & \textbf{Tests:} Students will be evaluated in test, descriptive
questions with different complexity levels (Revised Bloom's Taxonomy Levels:
Remembering, Understanding, Applying, Analyzing, Evaluating and Creating). TWO
tests will be conducted. Each test will be evaluated for 50 marks, adding up to
100 marks. Final test marks will be reduced to 40 marks. & \bfseries 40\\
\midrule
\rlab{3.} & \textbf{Experiential Learning:} Students will be evaluated for their
creativity and practical implementation of the problem. Case study-based
teaching learning (10), Program specific requirements (10), Video based
seminar/presentation/demonstration (20), adding up to 40. & \bfseries 40\\
\midrule
\multicolumn{2}{@{}l}{\bfseries MAXIMUM MARKS FOR THE CIE THEORY} & \bfseries 100\\
\bottomrule[1.1pt]
\end{longtable}}

\section*{Rubric for the Semester End Examination (Theory)}
\addcontentsline{toc}{section}{Rubric for the Semester End Examination}

{\sffamily\small
\begin{longtable}{@{}P{0.12\textwidth}P{0.72\textwidth}>{\raggedleft\arraybackslash}P{0.08\textwidth}@{}}
\toprule[1.1pt]
\rowcolor{ocre}\thd{Q. No.} & \thd{Contents} & \thd{Marks}\\
\midrule
\multicolumn{3}{@{}l}{\bfseries\color{slate} PART A}\\
\midrule
\rlab{1} & Objective type questions covering entire syllabus & \bfseries 20\\
\midrule
\multicolumn{3}{@{}l}{\bfseries\color{slate} PART B \quad\normalfont\itshape(Maximum of TWO sub-divisions only)}\\
\midrule
\rlab{2}      & Unit 1: (Compulsory)      & \bfseries 16\\
\rowcolor{tablealt}
\rlab{3 \& 4} & Unit 2: (Internal Choice) & \bfseries 16\\
\rlab{5 \& 6} & Unit 3: (Internal Choice) & \bfseries 16\\
\rowcolor{tablealt}
\rlab{7 \& 8} & Unit 4: (Internal Choice) & \bfseries 16\\
\rlab{9 \& 10}& Unit 5: (Internal Choice) & \bfseries 16\\
\midrule
\multicolumn{2}{@{}l}{\bfseries TOTAL} & \bfseries 100\\
\bottomrule[1.1pt]
\end{longtable}}

\vspace{8pt}

\begin{keybox}[What the SEE Structure Means for This Volume]
Part A draws twenty marks of objective questions from the \emph{entire}
syllabus, so no unit can be skipped. In Part B, Question~2 is compulsory and is
always set from Unit~I; Units~II to~V each carry a 16-mark internal-choice
pair. A candidate therefore answers Unit~I plus one question from each of the
remaining four units --- five full questions in all.
\end{keybox}


% ---------------------------------------------------------------- contents
\cleardoublepage
\pdfbookmark[0]{Contents}{toc}
\tableofcontents
\cleardoublepage
\listoffigures
\cleardoublepage
\listoftables

%%%%%%%%%%%%%%%%%%%%%%%%%%%%%%%%%%%%%%%%%%%%%%%%%%%%%%%%%%%%%%%%%%%%%%%%%%%%%%
\mainmatter
\pagestyle{mmmain}

%%%%%%%%%%%%%%%%%%%%%%%%%%%%%%%%%%%%%%%%%%%%%%%%%%%%%%%%%%%%%%%%%%%%%%%%%%%%%%
\chapter{AIDAS and Its Role}
%%%%%%%%%%%%%%%%%%%%%%%%%%%%%%%%%%%%%%%%%%%%%%%%%%%%%%%%%%%%%%%%%%%%%%%%%%%%%%

\syllabusstrip{Conversion Optimization: What is AIDAS and its role.}

\begin{learningoutcomes}
\begin{itemize}
\item What AIDAS\ix{AIDAS} stands for, and how it extends the classical AIDA model.
\item The five stages --- psychological state, marketer's task, and the website
      elements that deliver each.
\item Why Satisfaction is the critical component.
\item The role of AIDAS in conversion optimization --- as a diagnostic, not a
      description.
\item How to structure a product landing page with AIDAS.
\end{itemize}
\end{learningoutcomes}

\begin{keybox}[The Most Examined Topic in This Unit]
AIDAS is examined at both two and ten marks.
\pyq{SEE, Oct/Nov 2023; SEE, Jun/Jul 2025; Improvement CIE, 28 Aug 2024; Improvement CIE, 4 Jun 2025}
\end{keybox}

\section{What AIDAS Stands For}

\begin{exambox}[Model Answer --- 2 Marks]
\pyq{Improvement CIE, 4 Jun 2025 --- 2 M}
\textbf{A} --- Attention \quad\textbullet\quad \textbf{I} --- Interest
\quad\textbullet\quad \textbf{D} --- Desire \quad\textbullet\quad \textbf{A} ---
Action \quad\textbullet\quad \textbf{S} --- Satisfaction.
\end{exambox}

\keytermas{AIDAS}{AIDAS} is a hierarchy-of-effects model describing the sequence
of psychological states a prospect passes through, from first exposure to a brand
to post-purchase satisfaction. It extends the classical \keytermas{AIDA}{AIDA}
model by adding \textbf{Satisfaction}, which converts a one-time buyer into a
retained customer and an advocate.

\begin{figure}[htbp]
\centering
\begin{tikzpicture}[font=\sffamily\footnotesize]
  % narrowing bands, then a widening satisfaction band
  \fill[ocre!48] (-4.5,0)    -- (4.5,0)    -- (3.9,-1.05) -- (-3.9,-1.05) -- cycle;
  \fill[ocre!38] (-3.9,-1.05)-- (3.9,-1.05)-- (3.1,-2.10) -- (-3.1,-2.10) -- cycle;
  \fill[ocre!28] (-3.1,-2.10)-- (3.1,-2.10)-- (2.2,-3.15) -- (-2.2,-3.15) -- cycle;
  \fill[ocre!18] (-2.2,-3.15)-- (2.2,-3.15)-- (1.3,-4.20) -- (-1.3,-4.20) -- cycle;
  \fill[greenok!32] (-1.3,-4.20) -- (1.3,-4.20) -- (3.6,-5.35) -- (-3.6,-5.35) -- cycle;

  \node[text=slate,font=\sffamily\bfseries\small] at (0,-0.52) {ATTENTION};
  \node[text=slate,font=\sffamily\bfseries\small] at (0,-1.57) {INTEREST};
  \node[text=slate,font=\sffamily\bfseries\small] at (0,-2.62) {DESIRE};
  \node[text=slate,font=\sffamily\bfseries\small] at (0,-3.67) {ACTION};
  \node[text=greenok,font=\sffamily\bfseries\small] at (0,-4.80) {SATISFACTION};

  \node[anchor=west,text=slatelight,font=\sffamily\scriptsize\itshape] at (4.7,-0.52)
    {``something caught my eye''};
  \node[anchor=west,text=slatelight,font=\sffamily\scriptsize\itshape] at (4.7,-1.57)
    {``this might be relevant to me''};
  \node[anchor=west,text=slatelight,font=\sffamily\scriptsize\itshape] at (4.7,-2.62)
    {``I want this''};
  \node[anchor=west,text=slatelight,font=\sffamily\scriptsize\itshape] at (4.7,-3.67)
    {``I am buying it''};
  \node[anchor=west,text=greenok,font=\sffamily\scriptsize\itshape] at (4.7,-4.80)
    {``I am glad I did''};

  \draw[-{Latex[length=2.6mm]},greenok,line width=1pt]
    (-3.6,-5.35) .. controls (-6.0,-4.2) and (-6.0,-1.0) .. (-4.5,-0.15);
  \node[rotate=90,anchor=south,text=greenok,font=\sffamily\scriptsize\itshape]
    at (-5.85,-2.7) {advocacy returns as new attention};
\end{tikzpicture}
\caption[The AIDAS model]%
        {AIDAS. The first four stages narrow, as a funnel does. The fifth widens
         again, because a satisfied customer brings others --- which is precisely
         what the classical AIDA model, ending at Action, cannot express.}
\label{fig:aidas}
\end{figure}

\section{The Five Stages in Detail}

\begin{mmlongtable}{The five AIDAS stages}{tab:aidasstages}%
{@{}P{0.13\textwidth}P{0.20\textwidth}P{0.26\textwidth}P{0.30\textwidth}@{}}
\rowcolor{ocre}\thd{Stage} & \thd{Psychological state} & \thd{What the marketer must do} & \thd{Website elements that deliver it}\\
\midrule
\endfirsthead
\rowcolor{ocre}\thd{Stage} & \thd{Psychological state} & \thd{What the marketer must do} & \thd{Website elements that deliver it}\\
\midrule
\endhead
\rlab{Attention} & ``Something has caught my eye.'' & Break through indifference
in the first few seconds & Hero headline, striking visual, above-the-fold offer,
advertisement message match, page speed\\
\rowcolor{tablealt}
\rlab{Interest} & ``This might be relevant to me.'' & Convert curiosity into
engagement by showing relevance & Benefit-led sub-headline, scannable content,
product images and video, clear navigation\\
\rlab{Desire} & ``I want this.'' & Turn interest into preference through proof and
emotional appeal & Testimonials, ratings, case studies, comparison tables,
guarantees, scarcity and urgency cues, transparent pricing\\
\rowcolor{tablealt}
\rlab{Action} & ``I am buying it.'' & Remove every obstacle between wanting and
buying & One prominent call to action, short form, guest checkout, multiple
payment options, security badges, no hidden costs\\
\rlab{Satisfaction} & ``I am glad I did.'' & Ensure the delivered experience meets
or exceeds the promise & Order confirmation, proactive updates, easy returns,
responsive support, onboarding content, loyalty programme, review request\\
\end{mmlongtable}

\section{Why Satisfaction Is the Critical Component}

\begin{exambox}[Model Answer --- 2 Marks]
``Why is the `Satisfaction' component important in the AIDAS model?''
\pyq{SEE, Jun/Jul 2025 --- 2 M}

Because the sale is not the end of the funnel. Satisfaction determines whether
the customer returns, repurchases, recommends and defends the brand.

\smallskip
Its importance rests on the economics: the probability of selling to an existing
customer is \textbf{60--70\%} against \textbf{5--20\%} for a new prospect, and
repeat customers spend up to \textbf{67\% more}.

\smallskip
A dissatisfied customer, by contrast, does not merely fail to return --- in a
digital environment they broadcast the failure through reviews and social media,
damaging future acquisition. Satisfaction therefore converts a transaction into a
relationship, and turns customers into an unpaid distribution channel.
\end{exambox}

\section{The Role of AIDAS in Conversion Optimization}

\begin{keybox}[AIDAS Is a Diagnostic, Not a Description]
The model's value is that it tells you \emph{which element to fix}. Answering
``what is the role of AIDAS'' by simply re-listing the five stages misses the
question.
\pyq{SEE, Oct/Nov 2023 --- 2 M}
\end{keybox}

\begin{figure}[htbp]
\centering
\begin{tikzpicture}[font=\sffamily\footnotesize,
    sym/.style={draw=warnred,line width=0.85pt,fill=warnpale,rounded corners=3pt,
                minimum height=1.0cm,text width=3.5cm,align=center,
                font=\sffamily\scriptsize,text=slate,inner sep=4pt},
    dia/.style={draw=ocre,line width=0.85pt,fill=ocrepale,rounded corners=3pt,
                minimum height=1.0cm,text width=2.6cm,align=center,
                font=\sffamily\bfseries\scriptsize,text=ocredark,inner sep=4pt}]
  \node[anchor=south,text=warnred,font=\sffamily\bfseries\scriptsize] at (1.9,1.35)
    {THE SYMPTOM};
  \node[anchor=south,text=ocredark,font=\sffamily\bfseries\scriptsize] at (7.3,1.35)
    {THE STAGE THAT FAILED};

  \node[sym] (y1) at (1.9,0.75)   {high bounce, immediately};
  \node[sym] (y2) at (1.9,-0.45)  {high bounce \emph{after} scrolling};
  \node[sym] (y3) at (1.9,-1.65)  {high add-to-cart, low checkout};
  \node[sym] (y4) at (1.9,-2.85)  {high refunds, low repeat rate};

  \node[dia] (d1) at (7.3,0.75)   {ATTENTION};
  \node[dia] (d2) at (7.3,-0.45)  {INTEREST};
  \node[dia] (d3) at (7.3,-1.65)  {ACTION\\[-1pt]{\scriptsize\mdseries desire succeeded}};
  \node[dia,draw=greenok,fill=greenpale,text=greenok] (d4) at (7.3,-2.85)
     {SATISFACTION};

  \foreach \a/\b in {y1/d1,y2/d2,y3/d3,y4/d4}{
    \draw[-{Latex[length=2.2mm]},slate!60,line width=0.8pt] (\a)--(\b);}
\end{tikzpicture}
\caption[AIDAS as a diagnostic]%
        {Each characteristic symptom maps to exactly one failed stage. This is
         what makes AIDAS operationally useful rather than merely descriptive.}
\label{fig:aidasdiag}
\end{figure}

\begin{enumerate}
\item \sfb{It locates the failure.} High bounce\ix{bounce} means Attention failed. High
      bounce after scrolling means Interest failed. High add-to-cart with low
      checkout means Desire succeeded but Action failed. High refunds and a low
      repeat rate mean Satisfaction failed.
\item \sfb{It maps content to psychological state}, so each page section has a
      defined job rather than being decorative.
\item \sfb{It sequences the page} --- attention at the top, interest through the
      middle, desire immediately before the call to action\ix{call to action}, action at the point of
      decision, satisfaction promised in advance through guarantees and delivered
      afterwards.
\item \sfb{It generates testable hypotheses} --- ``the CTA fails because Desire is
      unsupported; adding ratings above the button will raise conversions.''
\item \sfb{It prevents over-optimising one stage.} A beautiful
      attention-grabbing page with no proof converts no better than a plain one.
\end{enumerate}

\section{Structuring a Product Landing Page With AIDAS}

\begin{keybox}[High Yield]
``Analyze how a company can use the AIDAS model to structure content on its
product landing page.''
\pyq{SEE, Jun/Jul 2025, Q10b --- 8 M}
\end{keybox}

\begin{figure}[htbp]
\centering
\begin{tikzpicture}[font=\sffamily\footnotesize,
    band/.style={minimum width=6.0cm,align=center,text width=5.7cm,
                 font=\sffamily\scriptsize,text=slate,inner sep=5pt,
                 anchor=north},
    tag/.style={anchor=west,font=\sffamily\bfseries\scriptsize}]

  % Bands are a uniform 1.85cm tall with a 0.10cm gutter, and each label is
  % trimmed to three lines, so no band's text can spill into the one below it.
  \draw[rulegrey,line width=0.9pt,rounded corners=3pt] (0,0.35) rectangle (6.2,-12.0);
  \fill[tablehead,rounded corners=3pt] (0,0.35) rectangle (6.2,-0.15);

  % --- attention
  \fill[ocre!30] (0.1,-0.25) rectangle (6.1,-2.10);
  \node[band] at (3.1,-0.36) {\textbf{HERO --- above the fold}\\[2pt]
     benefit headline matching the advertisement, product image or short video,
     offer visible without scrolling};
  \node[tag,text=ocredark] at (6.45,-1.18) {ATTENTION};

  % --- interest
  \fill[ocre!20] (0.1,-2.20) rectangle (6.1,-4.05);
  \node[band] at (3.1,-2.31) {\textbf{VALUE PROPOSITION BAND}\\[2pt]
     a sub-headline plus three or four benefit statements, in scannable form};
  \node[tag,text=ocredark] at (6.45,-3.13) {INTEREST};

  % --- desire (proof)
  \fill[ocre!13] (0.1,-4.15) rectangle (6.1,-6.00);
  \node[band] at (3.1,-4.26) {\textbf{PROOF SECTION}\\[2pt]
     ratings, testimonials with photographs, comparison table, awards and
     certifications};
  \node[tag,text=ocredark] at (6.45,-5.08) {DESIRE};

  % --- desire (offer)
  \fill[ocre!13] (0.1,-6.10) rectangle (6.1,-7.95);
  \node[band] at (3.1,-6.21) {\textbf{OFFER \& URGENCY BAND}\\[2pt]
     transparent pricing, inclusions, a countdown, a risk-reversal guarantee};
  \node[tag,text=ocredark] at (6.45,-7.03) {DESIRE};

  % --- action
  \fill[deepbluepale] (0.1,-8.05) rectangle (6.1,-9.90);
  \node[band] at (3.1,-8.16) {\textbf{CONVERSION BLOCK}\\[2pt]
     one dominant CTA, a short form, guest checkout, payment logos, a security
     badge};
  \node[tag,text=deepblue] at (6.45,-8.98) {ACTION};

  % --- satisfaction
  \fill[greenpale] (0.1,-10.00) rectangle (6.1,-11.85);
  \node[band] at (3.1,-10.11) {\textbf{REASSURANCE FOOTER}\\[2pt]
     delivery timelines, returns policy, support contact, FAQ};
  \node[tag,text=greenok] at (6.45,-10.93) {SATISFACTION};

  \draw[-{Latex[length=2.6mm]},slate!55,line width=1pt] (-0.55,0.15) -- (-0.55,-11.8);
  \node[rotate=90,anchor=south,text=slatelight,font=\sffamily\scriptsize\itshape]
    at (-0.75,-5.8) {scroll direction};
\end{tikzpicture}
\caption[A product landing page structured by AIDAS]%
        {Each band of the page carries exactly one psychological job, in the order
         the visitor experiences them.}
\label{fig:aidaspage}
\end{figure}

\begin{mmlongtable}{Page section mapped to AIDAS stage}{tab:aidaspage}%
{@{}P{0.22\textwidth}P{0.14\textwidth}P{0.57\textwidth}@{}}
\rowcolor{ocre}\thd{Page section} & \thd{AIDAS stage} & \thd{Content}\\
\midrule
\endfirsthead
\rowcolor{ocre}\thd{Page section} & \thd{AIDAS stage} & \thd{Content}\\
\midrule
\endhead
\rlab{Hero (above the fold)} & Attention & A benefit headline matching the
advertisement, a strong product image or short video, and the offer visible
without scrolling\\
\rowcolor{tablealt}
\rlab{Value proposition band} & Interest & A sub-headline plus three to four
benefit statements answering ``what's in it for me'', written in scannable form\\
\rlab{Proof section} & Desire & Star ratings, customer testimonials with
photographs, before-and-after visuals, a comparison table against alternatives,
awards and certifications\\
\rowcolor{tablealt}
\rlab{Offer and urgency band} & Desire & Transparent pricing, what is included, a
limited-time offer with a countdown, and a risk-reversal guarantee\\
\rlab{Conversion block} & Action & One dominant, high-contrast CTA button with
active benefit copy, a short form, guest checkout, payment logos and a security
badge\\
\rowcolor{tablealt}
\rlab{Reassurance footer} & Satisfaction & Delivery timelines, returns and
cancellation policy, support contact, FAQ, and a post-purchase onboarding
promise\\
\end{mmlongtable}

\begin{keybox}[The Structuring Principle]
Each section has \textbf{one} psychological job. If a section cannot be assigned
to an AIDAS stage, it is probably noise and should be deleted --- which is the
same conclusion the design principle ``omit needless words'' reaches from the
other direction.
\end{keybox}

\begin{examplebox}[One Example Per Stage]
\sfb{Attention} --- a streaming service's landing page whose entire first screen
is one headline, one CTA and nothing else.

\sfb{Interest} --- a product page opening with three benefit bullets rather than
technical specifications.

\sfb{Desire} --- a marketplace placing star ratings and review counts directly
above the buy button.

\sfb{Action} --- a one-tap checkout with saved payment details and guest checkout
enabled.

\sfb{Satisfaction} --- an order confirmation and proactive delivery-tracking
sequence followed by a review request, converting the buyer into an advocate.
\end{examplebox}

\begin{summarybox}
AIDAS is Attention, Interest, Desire, Action, Satisfaction --- a
hierarchy-of-effects model extending AIDA by adding the post-purchase stage that
turns a buyer into an advocate. Each stage has a psychological state, a marketer's
task and a set of website elements that deliver it. Satisfaction matters because
selling to an existing customer succeeds 60--70\% of the time against 5--20\% for
a new prospect, repeat customers spend up to 67\% more, and a dissatisfied
customer broadcasts the failure publicly. The model's role in conversion
optimization is diagnostic: it locates which stage failed from the symptom, maps
content to psychological state, sequences the page, generates testable hypotheses
and prevents over-optimising one stage. A landing page structured by AIDAS runs
hero, value proposition\ix{value proposition}, proof, offer and urgency, conversion block and
reassurance footer --- and any section that cannot be assigned a stage should be
deleted.
\end{summarybox}

%%%%%%%%%%%%%%%%%%%%%%%%%%%%%%%%%%%%%%%%%%%%%%%%%%%%%%%%%%%%%%%%%%%%%%%%%%%%%%
\chapter{Website Optimization for Conversion}
%%%%%%%%%%%%%%%%%%%%%%%%%%%%%%%%%%%%%%%%%%%%%%%%%%%%%%%%%%%%%%%%%%%%%%%%%%%%%%

\syllabusstrip{Conversion Optimization: website optimization.}

\begin{learningoutcomes}
\begin{itemize}
\item Conversion, conversion rate and conversion rate optimization defined.
\item Why CRO compounds, and why it is usually cheaper than buying traffic.
\item The five-stage conversion funnel\ix{conversion funnel}, its metrics and its typical drop-off
      causes.
\item How to identify drop-offs, and how to address them stage by stage.
\item The KPIs that evaluate website optimization, and what each contributes.
\item The seven-step optimisation workflow.
\end{itemize}
\end{learningoutcomes}

\section{Definitions}

\begin{definitionbox}[Conversion, Conversion Rate, and CRO]
A \keyterm{conversion} is the completion of a defined, valuable action by a
visitor --- a purchase, a form submission, a sign-up, a download, a call or an app
install.

\smallskip
The \keytermas{conversion rate}{conversion rate} is the percentage of visitors who
complete that action:
\[ \text{Conversion Rate} \;=\; \frac{\text{Conversions}}{\text{Total Visitors}}
   \times 100 \]

\keytermas{Conversion rate optimization}{conversion rate optimization} (CRO) is
the systematic process of increasing the percentage of visitors who complete a
desired action, by using analytics, user research and controlled experiments to
remove friction and strengthen persuasion.
\end{definitionbox}

\section{Why CRO Matters}

\begin{keybox}[The Compounding Argument]
\begin{itemize}
\item It \sfb{compounds on all future traffic} --- a 20\% lift applies to every
      visitor thereafter, at zero additional acquisition cost.
\item It is usually \sfb{cheaper than buying more traffic}; doubling the
      conversion rate has the same revenue effect as doubling the advertising
      budget, but without the recurring spend.
\item It \sfb{improves the ROI of every other channel simultaneously} --- SEO,
      paid, e-mail and social all deposit into the same funnel.
\item It \sfb{lowers customer acquisition cost}, and therefore improves the ratio
      of lifetime value to acquisition cost.
\item It \sfb{replaces opinion with evidence}, ending unresolvable internal
      debates about design.
\end{itemize}
\end{keybox}

\section{The Conversion Funnel}

\begin{keybox}[High Yield]
``Describe the conversion funnel in digital marketing. What are the key stages and
how can business identify and address drop offs in the funnel?''
\pyq{SEE, Oct/Nov 2023, Q9a --- 8 M}
The question has three parts --- stages, identification, remedy --- and all three
must be answered.
\end{keybox}

\begin{mmlongtable}{The five-stage conversion funnel}{tab:convfunnel}%
{@{}P{0.157\textwidth}P{0.235\textwidth}P{0.235\textwidth}P{0.285\textwidth}@{}}
\rowcolor{ocre}\thd{Stage} & \thd{What happens} & \thd{Typical metric} & \thd{Common cause of drop-off}\\
\midrule
\endfirsthead
\rowcolor{ocre}\thd{Stage} & \thd{What happens} & \thd{Typical metric} & \thd{Common cause of drop-off}\\
\midrule
\endhead
\rlab{Awareness} & The prospect first encounters the brand & Impressions, reach &
Wrong targeting; weak creative\\
\rowcolor{tablealt}
\rlab{Interest / Consideration} & They engage and evaluate & Sessions, pages per
session, time on site & Poor message match; unclear value proposition\\
\rlab{Desire / Intent} & They want the product and compare options & Add-to-cart
rate, product page views, wishlist adds & Price opacity; missing proof; weak
differentiation\\
\rowcolor{tablealt}
\rlab{Action / Conversion} & They complete the purchase or enquiry & Checkout
completion, form submission, conversion rate & Long forms; forced registration;
hidden costs; payment friction\\
\rlab{Satisfaction / Retention} & They receive value and return & Repeat purchase
rate, retention, net promoter score & Poor delivery; weak support; no
post-purchase communication\\
\end{mmlongtable}

\subsection{Identifying drop-offs}

\begin{enumerate}
\item \sfb{Build the funnel in analytics} --- define each step as an event and use
      funnel exploration to see the exact percentage lost at each transition.
\item \sfb{Segment the funnel} by device, channel, geography and new against
      returning; the aggregate figure usually hides where the real failure is.
\item \sfb{Add qualitative diagnosis} --- heatmaps, scroll maps and session
      recordings show \emph{why} users leave at the identified step.
\item \sfb{Ask the user directly} --- exit-intent surveys and post-purchase
      questions capture the reason in their own words.
\item \sfb{Compare against benchmarks} to judge whether a step is genuinely weak
      or simply normal for the category.
\end{enumerate}

\subsection{Addressing drop-offs}

\begin{itemize}
\item \sfb{Awareness \ra{} Interest:} improve message match\ix{message match} between the
      advertisement and the landing page; sharpen the headline.
\item \sfb{Interest \ra{} Desire:} add proof --- reviews, ratings,
      demonstrations, comparison tables; make pricing transparent.
\item \sfb{Desire \ra{} Action:} shorten the form, allow guest checkout, show
      total cost including shipping upfront, offer multiple payment methods, add a
      progress indicator, and place trust badges at the payment step.
\item \sfb{Action \ra{} Satisfaction:} confirm the order immediately, communicate
      proactively during fulfilment, make returns easy, and follow up for a
      review.
\item \sfb{Across all stages:} improve page speed, fix mobile-specific failures,
      and A/B test the single highest-loss step first.
\end{itemize}

\section{KPIs for Website Optimization}

\begin{keybox}[High Yield]
``What are the key performance indicators (KPIs) used to evaluate website
optimization, and how do they contribute to increased conversions?''
\pyq{SEE, Jun/Jul 2025, Q9a --- 8 M}
Note the second half of the question: each KPI must be paired with its
\emph{contribution}, not merely named.
\end{keybox}

\begin{mmlongtable}{KPIs for website optimization}{tab:croKPI}%
{@{}P{0.28\textwidth}P{0.65\textwidth}@{}}
\rowcolor{ocre}\thd{KPI} & \thd{Contribution to conversions}\\
\midrule
\endfirsthead
\rowcolor{ocre}\thd{KPI} & \thd{Contribution to conversions}\\
\midrule
\endhead
\rlab{Conversion rate} & The direct measure --- every other KPI exists to move
this one\\
\rowcolor{tablealt}
\rlab{Bounce rate} & Identifies pages where the visitor never enters the funnel at
all\\
\rlab{Average session duration and pages per session} & Indicate whether the
content sustains interest long enough to build desire\\
\rowcolor{tablealt}
\rlab{Page load time and Core Web Vitals} & Each second of delay removes a
measurable share of visitors before they see the offer\\
\rlab{Exit rate by page} & Pinpoints the specific page losing the customer\\
\rowcolor{tablealt}
\rlab{Cart abandonment rate} & Isolates friction in the final Action step, usually
the largest single recoverable loss\\
\rlab{Form completion rate} & Reveals whether the form itself is blocking the
conversion\\
\rowcolor{tablealt}
\rlab{Click-through rate on CTAs} & Tests whether the call to action is visible
and persuasive\\
\rlab{Scroll depth} & Determines what must be moved above the fold\\
\rowcolor{tablealt}
\rlab{Mobile against desktop conversion gap} & Exposes device-specific failures
hidden by the aggregate\\
\rlab{Return visitor rate and repeat purchase rate} & Measure the Satisfaction
stage and long-term value\\
\rowcolor{tablealt}
\rlab{Cost per acquisition and return on advertising spend} & Connect
optimisation work to profitability\\
\end{mmlongtable}

\section{The Optimization Workflow}

\begin{figure}[htbp]
\centering
\begin{tikzpicture}[font=\sffamily\footnotesize,
    st/.style={draw=ocre,line width=0.85pt,fill=ocrepale,rounded corners=3pt,
               minimum height=1.15cm,text width=1.62cm,align=center,
               font=\sffamily\bfseries\scriptsize,text=slate,inner sep=3pt}]
  \node[st] (w1) at (0,0)     {1\\ measure};
  \node[st] (w2) at (1.90,0)  {2\\ diagnose};
  \node[st] (w3) at (3.80,0)  {3\\ hypothesise};
  \node[st] (w4) at (5.70,0)  {4\\ prioritise};
  \node[st] (w5) at (7.60,0)  {5\\ test};
  \node[st] (w6) at (9.50,0)  {6\\ implement};
  \node[st,draw=greenok,fill=greenpale] (w7) at (11.40,0) {7\\ re-measure};
  \foreach \a/\b in {w1/w2,w2/w3,w3/w4,w4/w5,w5/w6,w6/w7}{
    \draw[-{Latex[length=1.9mm]},ocre,line width=0.85pt] (\a)--(\b);}
  \draw[-{Latex[length=2.1mm]},greenok,line width=0.85pt,dashed]
    (w7.south) -- ++(0,-0.7) -| (w1.south);
  \node[anchor=north,text=greenok,font=\sffamily\scriptsize\itshape]
    at (5.70,-0.75) {CRO is a continuous loop, not a project};
\end{tikzpicture}
\caption[The conversion optimisation workflow]%
        {The seven steps. Step~4 --- prioritising by expected impact against
         implementation effort --- is what stops the loop becoming an unordered
         backlog of ideas.}
\label{fig:croworkflow}
\end{figure}

\begin{enumerate}
\item \sfb{Measure} --- locate the weakest step in the funnel using analytics.
\item \sfb{Diagnose} --- use heatmaps, session recordings and surveys to find out
      why.
\item \sfb{Hypothesise} --- state a specific, falsifiable prediction tied to one
      variable.
\item \sfb{Prioritise} --- rank hypotheses by expected impact against
      implementation effort.
\item \sfb{Test} --- run an A/B or multivariate test to statistical significance.
\item \sfb{Implement} the winner.
\item \sfb{Re-measure and repeat} --- CRO is a continuous loop, not a project.
\end{enumerate}

\begin{summarybox}
A conversion is a defined valuable action; the conversion rate is conversions over
visitors, expressed as a percentage; CRO systematically raises that rate using
analytics, research and experiments. It matters because it compounds on all future
traffic at no extra acquisition cost, is usually cheaper than buying traffic,
improves every channel's ROI simultaneously, lowers acquisition cost and replaces
opinion with evidence. The funnel runs awareness, interest, desire, action and
satisfaction, each with its own metric and characteristic drop-off cause. Drop-offs
are identified by building the funnel in analytics, segmenting it, diagnosing
qualitatively, asking users directly and benchmarking --- and addressed stage by
stage from message match through proof and friction removal to post-purchase
communication. Twelve KPIs evaluate the work, each with a specific contribution.
The workflow is measure, diagnose, hypothesise, prioritise, test, implement,
re-measure --- and it loops.
\end{summarybox}

%%%%%%%%%%%%%%%%%%%%%%%%%%%%%%%%%%%%%%%%%%%%%%%%%%%%%%%%%%%%%%%%%%%%%%%%%%%%%%
\chapter{What Visitors Want to See on the Website}
%%%%%%%%%%%%%%%%%%%%%%%%%%%%%%%%%%%%%%%%%%%%%%%%%%%%%%%%%%%%%%%%%%%%%%%%%%%%%%

\syllabusstrip{Conversion Optimization: what visitors want to see on the website.}

\begin{learningoutcomes}
\begin{itemize}
\item The ten questions a visitor arrives with, and what must be present to
      answer each.
\item The psychology beneath them --- Krug's first law, scanning and satisficing.
\end{itemize}
\end{learningoutcomes}

\begin{keybox}[Examined at Both Weights]
This is asked as a two-mark question --- ``two key elements visitors expect on a
well-optimized landing page'' --- and as a ten-mark question requiring the full
set with delivery strategies.
\pyq{Improvement CIE, 4 Jun 2025 --- 2 M and 10 M}
\end{keybox}

\section{The Ten Visitor Questions}

\begin{mmlongtable}{What visitors want, and what must be present}{tab:visitorwants5}%
{@{}P{0.32\textwidth}P{0.61\textwidth}@{}}
\rowcolor{ocre}\thd{Visitor question} & \thd{What must be present}\\
\midrule
\endfirsthead
\rowcolor{ocre}\thd{Visitor question} & \thd{What must be present}\\
\midrule
\endhead
\rlab{``Am I in the right place?''} & A clear headline stating what is offered and
for whom, above the fold\\
\rowcolor{tablealt}
\rlab{``What's in it for me?''} & Benefit-led copy rather than feature lists or
company history\\
\rlab{``Can I trust you?''} & Testimonials, star ratings, client logos,
certifications, guarantees, real photographs, a physical address\\
\rowcolor{tablealt}
\rlab{``What does it cost?''} & Transparent pricing or at least a clear range;
hidden pricing is a leading cause of abandonment\\
\rlab{``What do I do next?''} & One prominent, unambiguous call to action\\
\rowcolor{tablealt}
\rlab{``Can I find what I came for?''} & Simple, conventional navigation and a
working search\\
\rlab{``Is it fast?''} & A page loading in under about three seconds\\
\rowcolor{tablealt}
\rlab{``Does it work on my phone?''} & A fully responsive, mobile-first layout\\
\rlab{``Is it safe to pay here?''} & HTTPS, recognised payment logos, a clear
privacy and returns policy\\
\rowcolor{tablealt}
\rlab{``How do I reach a human?''} & Visible contact options and stated response
times\\
\end{mmlongtable}

\section{The Psychology Behind It}

\begin{keybox}[Krug's First Law --- Don't Make Me Think]
A person of average or below-average ability must be able to work out how to
accomplish the task. It must be self-evident and self-explanatory.
\end{keybox}

Users do not read pages; they \keyterm{scan}, \keyterm{satisfice} and muddle
through. They choose the first option that plausibly works rather than the best
one. Three reasons they do not read instructions: it isn't important to them; they
are on a mission and know they don't need everything; and if something works they
stick with it rather than seeking a better way.

Design must therefore follow \sfb{visual hierarchy}, \sfb{omit needless words},
\sfb{make clickables obvious}, and \sfb{provide breadcrumbs and street signs}.

\begin{warnbox}[The Link Between This Chapter and AIDAS]
The ten questions are not an arbitrary checklist. Read them against
Chapter~1 and they resolve into the model: questions one and two are
\emph{Attention} and \emph{Interest}; trust, cost and safety are \emph{Desire};
``what do I do next'' is \emph{Action}; and ``how do I reach a human'' is the
promise of \emph{Satisfaction} made in advance. Answering the ten questions
\emph{is} building an AIDAS-complete page.
\end{warnbox}

\begin{summarybox}
Visitors arrive with ten questions: am I in the right place, what's in it for me,
can I trust you, what does it cost, what do I do next, can I find what I came for,
is it fast, does it work on my phone, is it safe to pay, and how do I reach a
human. Each has a specific delivery mechanism. Beneath them lies Krug's first law
--- don't make me think --- and the fact that users scan, satisfice and muddle
through rather than reading, which is why pages must follow visual hierarchy\ix{visual hierarchy}, omit
needless words, make clickables obvious and provide clear navigation. The ten
questions map onto the AIDAS stages, so answering them all is the same thing as
building an AIDAS-complete page.
\end{summarybox}

%%%%%%%%%%%%%%%%%%%%%%%%%%%%%%%%%%%%%%%%%%%%%%%%%%%%%%%%%%%%%%%%%%%%%%%%%%%%%%
\chapter{Optimizing Key Elements of a Landing Page}
%%%%%%%%%%%%%%%%%%%%%%%%%%%%%%%%%%%%%%%%%%%%%%%%%%%%%%%%%%%%%%%%%%%%%%%%%%%%%%

\syllabusstrip{Conversion Optimization: how to optimize key element and increase
the effect of landing on a particular page.}

\begin{learningoutcomes}
\begin{itemize}
\item The headline --- message match, specificity and dominance.
\item The call to action --- dominance, contrast, copy, position, size and
      reassurance.
\item Visuals --- authenticity, faces, video and direction.
\item Forms --- field count, guest checkout, smart\ix{SMART} defaults and inline
      validation.
\item Trust signals and social proof --- the concept, the impact and the ten
      types.
\item Page speed and mobile as conversion elements.
\end{itemize}
\end{learningoutcomes}

\section{The Headline}

\begin{itemize}
\item \sfb{Restate the promise} made by the advertisement, e-mail or search result
      that brought the visitor --- \keytermas{message match}{message match} is the
      single largest lever on landing-page conversion.
\item \sfb{Be specific and quantified} --- ``Kerala backwaters, 4 nights from
      \rupee 18,999'' beats ``Unforgettable holidays''.
\item \sfb{Lead with the benefit}, not the product name or the company name.
\item \sfb{Keep it short} enough to be read in one glance, and set it as the
      visually dominant element on the page.
\item \sfb{Support it with a sub-headline} that adds the second most important
      piece of information.
\end{itemize}

\section{The Call to Action}

\begin{itemize}
\item \sfb{One dominant CTA per page.} Two competing CTAs halve the conversion
      rate of both.
\item \sfb{High contrast} --- reserve one accent colour exclusively for actions,
      so that the eye finds it instantly.
\item \sfb{Active, benefit-led, first-person copy} --- ``Get my free trial''
      outperforms ``Submit'' or ``Click here''.
\item \sfb{Positioned above the fold} and repeated at natural decision points down
      the page.
\item \sfb{Adequate size and tap target} --- at least 44 by 44 pixels on mobile.
\item \sfb{Reassurance microcopy beneath the button} --- ``No card required
      \textbullet\ Cancel anytime \textbullet\ Takes 30 seconds''.
\end{itemize}

\section{Visuals}

\begin{itemize}
\item \sfb{Show the actual product or outcome}, not generic stock imagery, which
      reads as inauthentic and lowers trust.
\item \sfb{Use real customer photographs} in testimonials --- faces materially
      increase credibility.
\item \sfb{Add short video} --- 77\% of consumers have been convinced to buy after
      watching a video.
\item \sfb{Direct attention with the image} --- a person in the image looking
      toward the CTA measurably increases attention to it.
\item \sfb{Optimise every image} for weight and format, so that the visual does
      not destroy the speed it was meant to support.
\end{itemize}

\section{Forms}

\begin{itemize}
\item \sfb{Remove every field you will not act on.} Each additional field reduces
      completion.
\item \sfb{Allow guest checkout} --- forced registration is one of the largest
      single causes of cart abandonment.
\item \sfb{Use smart defaults, auto-fill and postcode lookup} to reduce typing on
      mobile.
\item \sfb{Validate inline, not on submit}, and write error messages that say how
      to fix the problem.
\item \sfb{Show a progress indicator} for multi-step forms.
\item \sfb{Use the right keyboard type} on mobile for each field.
\end{itemize}

\section{Trust Signals and Social Proof}

\begin{keybox}[High Yield]
``Explain the concept of social proof and its impact on conversion rates. Provide
examples of how business can use social proof effectively in marketing.''
\pyq{SEE, Oct/Nov 2023, Q9b --- 8 M}
\end{keybox}

\begin{definitionbox}[Social Proof]
\keytermas{Social proof}{social proof} is the psychological phenomenon whereby
people assume the actions and opinions of others reflect the correct behaviour in
a given situation --- so, uncertain about a decision, a buyer looks at what other
buyers did.
\end{definitionbox}

\begin{exambox}[Impact on Conversion Rates]
Social proof reduces \textbf{perceived risk}, which is the dominant barrier at the
Desire~\ra~Action transition. It substitutes borrowed credibility for a brand the
visitor does not yet know: roughly 90\% of people believe what peers say about a
brand while only about 48\% trust what businesses say directly. It also shortens
deliberation by supplying a decision shortcut, and it raises average order value
when the proof is attached to higher-tier options.
\end{exambox}

\begin{mmlongtable}{The ten types of social proof}{tab:socialproof}%
{@{}P{0.26\textwidth}P{0.67\textwidth}@{}}
\rowcolor{ocre}\thd{Type of social proof} & \thd{How to use it}\\
\midrule
\endfirsthead
\rowcolor{ocre}\thd{Type of social proof} & \thd{How to use it}\\
\midrule
\endhead
\rlab{Customer testimonials} & Short quotes with a real name, photograph and
specific outcome, placed adjacent to the call to action\\
\rowcolor{tablealt}
\rlab{Star ratings and reviews} & Aggregate rating on product pages and in search
rich results; display negative reviews too, since an all-perfect page reads as
fake\\
\rlab{User counts} & ``Joined by 40,000 students''; ``2 million orders
delivered''\\
\rowcolor{tablealt}
\rlab{Case studies} & Detailed, quantified success stories, for
high-consideration B2B purchases\\
\rlab{Client and media logos} & ``As used by'' and ``As featured in'' bands\\
\rowcolor{tablealt}
\rlab{Expert endorsement} & Certifications, professional bodies, specialist
reviews\\
\rlab{Influencer and celebrity proof} & Borrowed audience trust, disclosed as
sponsored\\
\rowcolor{tablealt}
\rlab{User-generated content} & Real customer photographs and videos, reposted
with permission\\
\rlab{Real-time activity notifications} & ``12 people are viewing this now'';
``Booked 4 times today'' --- effective, but must be truthful\\
\rowcolor{tablealt}
\rlab{Wisdom of the crowd} & ``Most popular plan'' and ``Best seller'' badges
guiding choice\\
\end{mmlongtable}

\begin{warnbox}[Two Ways Social Proof Backfires]
An \textbf{all-perfect review page reads as fabricated}; displaying some negative
reviews raises credibility rather than lowering it. And \textbf{real-time activity
notifications must be true} --- a fabricated ``12 people viewing'' is both a
consumer-protection risk and, once noticed, more damaging to trust than showing
nothing at all.
\end{warnbox}

\section{Page Speed and Mobile}

\begin{itemize}
\item A page loading in over three seconds loses a large share of mobile visitors
      before it is ever seen --- \sfb{speed is a conversion element, not a
      technical detail}.
\item Compress and modernise images, minify scripts, enable caching and a content
      delivery network, and lazy-load below-the-fold content.
\item Design mobile-first, since most traffic and most abandonment now occur on
      mobile.
\end{itemize}

\begin{summarybox}
The headline must restate the advertisement's promise, be specific and quantified,
lead with the benefit and dominate the page. The call to action must be single,
high-contrast, actively worded, above the fold and repeated, large enough to tap,
and supported by reassurance microcopy. Visuals must be authentic, feature real
faces, include short video, direct attention toward the CTA and be optimised for
weight. Forms must drop unused fields, allow guest checkout, use smart defaults,
validate inline and show progress. Social proof reduces perceived risk at the
Desire-to-Action transition, substituting borrowed credibility in ten forms from
testimonials to wisdom-of-the-crowd badges --- but it must be honest, since
fabricated proof destroys more trust than it borrows. Speed and mobile-first design
are conversion elements in their own right.
\end{summarybox}

%%%%%%%%%%%%%%%%%%%%%%%%%%%%%%%%%%%%%%%%%%%%%%%%%%%%%%%%%%%%%%%%%%%%%%%%%%%%%%
%  Unit V, second half: Digital Analytics (Chapters 5-7)
%%%%%%%%%%%%%%%%%%%%%%%%%%%%%%%%%%%%%%%%%%%%%%%%%%%%%%%%%%%%%%%%%%%%%%%%%%%%%%

\chapter{The Evolution of Digital Analytics}

\syllabusstrip{Digital Analytics: Evolution of Digital Analytics.}

\begin{learningoutcomes}
\begin{itemize}
\item What digital analytics is, and how it differs from web analytics.
\item How digital analytics influences business decisions.
\item The six eras of the discipline --- what each measured, with what tools, and
      what each could not do.
\item The direction of travel across all six.
\end{itemize}
\end{learningoutcomes}

\section{Definition}

\begin{definitionbox}[Digital Analytics]
\keytermas{Digital analytics}{digital analytics} is the collection, measurement,
analysis and interpretation of digital data in order to improve business
performance and customer experience.
\pyq{CIE-III, Jun 2026 --- 2 M}
\end{definitionbox}

\begin{keybox}[Digital Analytics Is Not Web Analytics]
The distinction matters and is worth stating explicitly in any answer.

\smallskip
\textbf{Web analytics} is confined to website behaviour. \textbf{Digital
analytics} is broader --- it unifies website, mobile app, social, e-mail, paid
media, CRM, offline and support data into a single view of the customer, and its
purpose is to \emph{inform business decisions} rather than merely to report on a
channel.
\end{keybox}

\section{How Digital Analytics Influences Business Decisions}

\begin{exambox}[Model Answer --- 2 Marks]
\pyq{CIE-III, Jun 2026 --- 2 M}
Digital analytics provides customer insights and performance data that help
businesses improve marketing strategies and customer engagement.

\smallskip
Concretely, it determines where budget is allocated, which products are
developed, which segments are targeted, which experiences are fixed first, and
what the business forecasts.
\end{exambox}

\section{The Six Eras}

\begin{keybox}[High Yield]
``Trace the evolution of digital analytics and its growing importance in
delivering a seamless end-to-end customer experience.''
\pyq{Improvement CIE, 4 Jun 2025 --- 10 M}
Give the eras \emph{with their limitations}: the limitation of each era is what
motivated the next, and that causal chain is the argument the question is asking
for.
\end{keybox}

\begin{figure}[htbp]
\centering
\begin{tikzpicture}[font=\sffamily\footnotesize,
    era/.style={draw=ocre,line width=0.85pt,fill=ocrepale,rounded corners=3pt,
                minimum height=1.5cm,text width=1.55cm,align=center,
                font=\sffamily\bfseries\scriptsize,text=slate,inner sep=3pt}]
  \draw[rulegrey,line width=1pt] (-0.4,0) -- (11.6,0);

  \node[era] (e1) at (0.55,1.15)  {1\\ server logs\\[-1pt]{\scriptsize\mdseries 1990s}};
  \node[era] (e2) at (2.75,-1.15) {2\\ page-tag web analytics\\[-1pt]{\scriptsize\mdseries early 2000s}};
  \node[era] (e3) at (4.75,1.15)  {3\\ visitor \& goal analytics\\[-1pt]{\scriptsize\mdseries mid 2000s}};
  \node[era] (e4) at (6.75,-1.15) {4\\ multi-channel \& social\\[-1pt]{\scriptsize\mdseries 2010s}};
  \node[era] (e5) at (8.75,1.15)  {5\\ user-centric, event-based\\[-1pt]{\scriptsize\mdseries late 2010s}};
  \node[era,draw=greenok,fill=greenpale] (e6) at (10.85,-1.15)
     {6\\ unified, predictive, privacy-first\\[-1pt]{\scriptsize\mdseries present}};

  \foreach \n/\x in {e1/0.55,e3/4.75,e5/8.75}{
    \draw[ocre,line width=0.8pt] (\x,0) -- (\x,0.4);
    \fill[ocre] (\x,0) circle (0.075);}
  \foreach \n/\x in {e2/2.75,e4/6.75}{
    \draw[ocre,line width=0.8pt] (\x,0) -- (\x,-0.4);
    \fill[ocre] (\x,0) circle (0.075);}
  \draw[greenok,line width=0.8pt] (10.85,0) -- (10.85,-0.4);
  \fill[greenok] (10.85,0) circle (0.075);

  \node[anchor=north,text=slatelight,font=\sffamily\scriptsize\itshape,
        align=center,text width=11.8cm] at (5.6,-2.35)
    {counting files \ra{} counting visits \ra{} counting conversions \ra{}
     understanding people \ra{} predicting behaviour};
\end{tikzpicture}
\caption[The six eras of digital analytics]%
        {Six eras. Each was created by the limitation of the one before it, and
         each moved analytics closer to the customer and closer to the boardroom.}
\label{fig:sixeras}
\end{figure}

\begin{mmlongtable}{The six eras of digital analytics}{tab:sixeras}%
{@{}P{0.167\textwidth}P{0.245\textwidth}P{0.235\textwidth}P{0.265\textwidth}@{}}
\rowcolor{ocre}\thd{Era} & \thd{What was measured} & \thd{Tools and methods} & \thd{Limitation}\\
\midrule
\endfirsthead
\rowcolor{ocre}\thd{Era} & \thd{What was measured} & \thd{Tools and methods} & \thd{Limitation}\\
\midrule
\endhead
\rlab{1. Server log analysis} \newline {\scriptsize 1990s} & Hits, file requests,
bandwidth & Raw server logs, log parsers & Counted files, not people; caches and
bots corrupted the data\\
\rowcolor{tablealt}
\rlab{2. Page-tag web analytics} \newline {\scriptsize early 2000s} & Pageviews,
visits, referrers, bounce rate & JavaScript tags; the first hosted analytics
platforms & Session-based and site-siloed; no view of the person across visits\\
\rlab{3. Visitor and goal analytics} \newline {\scriptsize mid to late 2000s} &
Goals, funnels, e-commerce revenue, segments & Goal tracking, funnel
visualisation, A/B testing tools & Last-click attribution; ignored the rest of the
journey\\
\rowcolor{tablealt}
\rlab{4. Multi-channel and social analytics} \newline {\scriptsize 2010s} &
Cross-channel attribution, social listening, mobile app analytics & Multi-channel
funnels, mobile SDKs, social insight platforms & Data still fragmented across
separate platforms\\
\rlab{5. User-centric and event-based analytics} \newline {\scriptsize late 2010s
onward} & The individual user across devices and sessions; every interaction as an
event & Event-based analytics platforms, tag managers, cross-device identity &
Requires disciplined measurement planning and governance\\
\rowcolor{tablealt}
\rlab{6. Unified, predictive and privacy-first analytics} \newline {\scriptsize
present} & End-to-end customer experience; predicted churn and lifetime value;
consent-governed first-party data & Customer data platforms, data warehouses,
machine learning, server-side tagging, consent management & Privacy regulation and
cookie deprecation constrain tracking; modelling replaces some observation\\
\end{mmlongtable}

\begin{keybox}[The Direction of Travel]
The movement is consistent across all six eras: from \textbf{counting files
\ra{} counting visits \ra{} counting conversions \ra{} understanding people
\ra{} predicting behaviour}, and from \textbf{channel-siloed reporting \ra{} a
unified end-to-end view}.

\smallskip
Each shift moved analytics closer to the customer and closer to the boardroom.
\end{keybox}

\begin{summarybox}
Digital analytics is the collection, measurement, analysis and interpretation of
digital data to improve business performance and customer experience --- broader
than web analytics because it unifies web, app, social, e-mail, paid, CRM, offline
and support data, and because its purpose is decision rather than report. It
influences budget allocation, product development, segment targeting, experience
priorities and forecasting. Six eras mark its evolution: server logs, page-tag web
analytics, visitor and goal analytics, multi-channel and social analytics,
user-centric event-based analytics, and unified predictive privacy-first
analytics. Each era's limitation created the next, and the direction of travel runs
from counting files to predicting behaviour, and from silos to a unified view.
\end{summarybox}

%%%%%%%%%%%%%%%%%%%%%%%%%%%%%%%%%%%%%%%%%%%%%%%%%%%%%%%%%%%%%%%%%%%%%%%%%%%%%%
\chapter{The End-to-End Customer Experience}
%%%%%%%%%%%%%%%%%%%%%%%%%%%%%%%%%%%%%%%%%%%%%%%%%%%%%%%%%%%%%%%%%%%%%%%%%%%%%%

\syllabusstrip{Digital Analytics: information about end-to-end customer
experience.}

\begin{learningoutcomes}
\begin{itemize}
\item What the end-to-end customer experience means in digital analytics.
\item Why the siloed, channel-by-channel view fails.
\item The seven requirements for genuine end-to-end measurement.
\item The business value it produces.
\item How market segmentation is refined using digital analytics data.
\end{itemize}
\end{learningoutcomes}

\section{What It Means}

\begin{exambox}[Model Answer --- 2 Marks]
``What is meant by the `end-to-end customer experience' in digital analytics?''
\pyq{SEE, Jun/Jul 2025 --- 2 M}

The end-to-end customer experience means measuring and understanding the
\textbf{complete customer journey\ix{customer journey} across every touchpoint and every stage} ---
from the first anonymous impression\ix{impression}, through research, comparison, purchase,
delivery and support, to repeat purchase and advocacy\ix{advocacy} --- as a \textbf{single
connected experience} rather than as a set of disconnected channel reports.
\end{exambox}

\section{Why the Siloed View Fails}

\begin{itemize}
\item A visitor may discover the brand on one social platform, research it on
      video, compare it on search, read reviews on a marketplace and buy in a
      store. Channel-level reporting credits only the last touch, and therefore
      misleads budget allocation.
\item The modern journey is not a funnel. It is a \keytermas{decision
      constellation}{decision constellation} --- a network of micro-decisions
      happening in different places at the same time. Optimising for visibility in
      one place while the customer is deciding everywhere else produces the classic
      symptom: traffic looks healthy but conversions are flat.
\item Silos also hide experience failures that span teams --- for example, a
      marketing promise that the fulfilment process cannot keep.
\end{itemize}

\begin{figure}[htbp]
\centering
\begin{tikzpicture}[font=\sffamily\footnotesize]
  % --- left: the siloed view
  \node[anchor=north,text=warnred,font=\sffamily\bfseries\scriptsize] at (2.3,2.5)
    {THE SILOED VIEW};
  \foreach \i/\x in {1/0.6,2/1.75,3/2.9,4/4.05}{
    \fill[warnpale] (\x-0.42,-1.5) rectangle (\x+0.42,1.9);
    \draw[warnred!60,line width=0.8pt] (\x-0.42,-1.5) rectangle (\x+0.42,1.9);}
  \node[rotate=90,text=slate,font=\sffamily\scriptsize] at (0.6,0.2)  {search};
  \node[rotate=90,text=slate,font=\sffamily\scriptsize] at (1.75,0.2) {social};
  \node[rotate=90,text=slate,font=\sffamily\scriptsize] at (2.9,0.2)  {e-mail};
  \node[rotate=90,text=slate,font=\sffamily\scriptsize] at (4.05,0.2) {store};
  \node[anchor=north,text=slatelight,font=\sffamily\scriptsize\itshape,
        align=center,text width=4.4cm] at (2.3,-1.65)
    {four reports, four last-click winners, one customer nobody can see};

  % --- right: the unified view
  \node[anchor=north,text=greenok,font=\sffamily\bfseries\scriptsize] at (8.9,2.5)
    {THE UNIFIED VIEW};
  \draw[greenok,line width=1pt,rounded corners=5pt,fill=greenpale]
    (6.5,-1.5) rectangle (11.3,1.9);
  \foreach \a/\lbl in {150/search,90/social,30/video,-30/reviews,-90/store,-150/support}{
    \node[draw=greenok!70,line width=0.7pt,fill=white,rounded corners=2pt,
          inner sep=2.2pt,font=\sffamily\scriptsize,text=slate]
      (t\lbl) at ($(8.9,0.2)+(\a:1.72cm and 1.05cm)$) {\lbl};
    \draw[greenok!60,line width=0.6pt,dashed] (t\lbl) -- (8.9,0.2);}
  \node[circle,fill=greenok,text=white,inner sep=2.2pt,
        font=\sffamily\bfseries\scriptsize] at (8.9,0.2) {1};
  \node[anchor=north,text=slatelight,font=\sffamily\scriptsize\itshape,
        align=center,text width=4.8cm] at (8.9,-1.65)
    {one identity, one event taxonomy, multi-touch credit --- one journey};

  \draw[-{Latex[length=2.6mm]},slate,line width=1pt] (4.75,0.2) -- (6.35,0.2);
\end{tikzpicture}
\caption[The siloed view against the unified view]%
        {Four channel reports do not add up to one customer. The unified view is
         what makes multi-touch attribution, journey analysis and
         cross-channel personalisation possible at all.}
\label{fig:siloedunified}
\end{figure}

\section{What End-to-End Analytics Requires}

\begin{enumerate}
\item A \sfb{unified customer identity} across web, app, e-mail, CRM and offline
      systems.
\item A \sfb{consistent event taxonomy}, so that the same action means the same
      thing in every system.
\item \sfb{Journey and path analysis}, not just page reports.
\item \sfb{Multi-touch attribution} rather than last-click.
\item \sfb{Voice-of-customer data} --- surveys, net promoter score, reviews and
      support tickets --- joined to behavioural data.
\item \sfb{Post-purchase and service data}, so that the measurement does not stop
      at the transaction.
\item \sfb{Consent and privacy governance} built into the collection layer.
\end{enumerate}

\section{Business Value}

\begin{itemize}
\item \sfb{Accurate budget allocation}, based on true contribution rather than
      last-click credit.
\item \sfb{Identification of the highest-friction moment} in the whole journey,
      not merely on one page.
\item \sfb{Personalisation that is consistent} across channels rather than
      contradictory.
\item \sfb{Improved retention}, because the analysis continues past the purchase.
\item \sfb{Predictive capability} --- churn\ix{churn} risk and lifetime value modelled from
      journey behaviour.
\end{itemize}

\section{Refining Market Segmentation With Digital Analytics}

\begin{keybox}[High Yield]
``Analyze how market segmentation can be refined using digital analytics data.''
\pyq{SEE, Jun/Jul 2025, Q9b --- 8 M}
\end{keybox}

\begin{itemize}
\item \sfb{From demographic to behavioural segmentation} --- traditional
      segmentation groups by age and income; analytics groups by what people
      actually do, which predicts purchase far better.
\item \sfb{Value-based segmentation} --- \keytermas{RFM analysis}{RFM analysis}
      (recency, frequency, monetary value) and lifetime-value modelling identify
      the segments worth investing in.
\item \sfb{Journey-stage segmentation} --- separating first-time visitors,
      researchers, cart abandoners, first-time buyers and loyalists, each
      requiring different messaging.
\item \sfb{Intent-based segmentation} --- search query and on-site behaviour
      reveal purchase intent in real time.
\item \sfb{Channel-affinity segmentation} --- which segment responds to which
      channel, enabling efficient allocation.
\item \sfb{Predictive segmentation} --- machine learning models group users by
      predicted churn risk or predicted conversion probability rather than by past
      attributes.
\item \sfb{Continuous refinement} --- analytics-driven segments update
      automatically as behaviour changes, whereas survey-based segments decay from
      the day they are created.
\end{itemize}

\begin{summarybox}
The end-to-end customer experience means measuring the complete journey across
every touchpoint and stage as one connected experience rather than as disconnected
channel reports. The siloed view fails because last-click credit misleads budget
allocation, because the modern journey is a decision constellation rather than a
funnel, and because silos hide failures that span teams. Genuine end-to-end
measurement requires unified identity, a consistent event taxonomy, journey
analysis, multi-touch attribution, voice-of-customer data, post-purchase data and
privacy governance. Its value is accurate budget allocation, identification of the
worst friction in the whole journey, consistent personalisation, improved
retention and predictive capability. Analytics refines segmentation from
demographic to behavioural, value-based, journey-stage, intent-based,
channel-affinity and predictive --- and keeps those segments current
automatically.
\end{summarybox}

%%%%%%%%%%%%%%%%%%%%%%%%%%%%%%%%%%%%%%%%%%%%%%%%%%%%%%%%%%%%%%%%%%%%%%%%%%%%%%
\chapter{The Analyst's Influence and Role as a Change Agent}
%%%%%%%%%%%%%%%%%%%%%%%%%%%%%%%%%%%%%%%%%%%%%%%%%%%%%%%%%%%%%%%%%%%%%%%%%%%%%%

\syllabusstrip{Digital Analytics: analyst's influence on business, role as a
change agent.}

\begin{learningoutcomes}
\begin{itemize}
\item Why data requires interpretation, and where the chain can break.
\item The seven ways an analyst influences the business.
\item What separates a reporting function from a change agent.
\item What being a change agent requires, and the barriers it must overcome.
\item Designing a holistic marketing\ix{holistic marketing} dashboard.
\item Applying digital analytics to two worked business problems.
\end{itemize}
\end{learningoutcomes}

\section{Data Requires Interpretation}

A theme repeated throughout this unit: data requires interpretation, and errors
compound at every step of the chain.

\begin{figure}[htbp]
\centering
\begin{tikzpicture}[font=\sffamily\footnotesize,
    stp/.style={draw=slate!55,line width=0.85pt,fill=white,rounded corners=3pt,
                minimum height=1.05cm,text width=1.85cm,align=center,
                font=\sffamily\bfseries\scriptsize,text=slate,inner sep=3pt}]
  \node[stp] (c1) at (0,0)     {1. product design};
  \node[stp,fill=ocrepale,draw=ocre] (c2) at (2.35,0)  {2. tool setup};
  \node[stp] (c3) at (4.70,0)  {3. data output};
  \node[stp,fill=ocrepale,draw=ocre] (c4) at (7.05,0)  {4. human interpretation};
  \node[stp,fill=ocrepale,draw=ocre] (c5) at (9.40,0)  {5. business action};
  \foreach \a/\b in {c1/c2,c2/c3,c3/c4,c4/c5}{
    \draw[-{Latex[length=1.9mm]},slate!60,line width=0.85pt] (\a)--(\b);}

  \node[anchor=north,align=center,text width=2.15cm,font=\sffamily\scriptsize,
        text=warnred] at (0,-0.62)    {the product may be flawed};
  \node[anchor=north,align=center,text width=2.15cm,font=\sffamily\scriptsize,
        text=warnred] at (2.35,-0.62) {mis-tagged pages, duplicate tracking,
        wrong filters};
  \node[anchor=north,align=center,text width=2.15cm,font=\sffamily\scriptsize,
        text=warnred] at (4.70,-0.62) {sampling, bot traffic, consent gaps};
  \node[anchor=north,align=center,text width=2.15cm,font=\sffamily\scriptsize,
        text=warnred] at (7.05,-0.62) {correlation read as causation; averages
        hiding segments};
  \node[anchor=north,align=center,text width=2.15cm,font=\sffamily\scriptsize,
        text=warnred] at (9.40,-0.62) {the final decision may be wrong};

  \draw[decorate,decoration={brace,amplitude=5pt,mirror},ocre,line width=0.8pt]
    (1.35,-2.55) -- (10.4,-2.55);
  \node[anchor=north,text=ocredark,font=\sffamily\bfseries\scriptsize]
    at (5.9,-2.75) {the analyst's value lies in steps 2, 4 and 5};
\end{tikzpicture}
\caption[The interpretation chain]%
        {Five steps, each of which can break. The analyst does not control step~1
         and only partly controls step~3, which is why their contribution is
         concentrated in instrumenting, interpreting and translating.}
\label{fig:interpchain}
\end{figure}

\begin{mmlongtable}{Where the interpretation chain breaks}{tab:interpchain}%
{@{}P{0.19\textwidth}P{0.35\textwidth}P{0.39\textwidth}@{}}
\rowcolor{ocre}\thd{Step} & \thd{What happens} & \thd{Where it can go wrong}\\
\midrule
\endfirsthead
\rowcolor{ocre}\thd{Step} & \thd{What happens} & \thd{Where it can go wrong}\\
\midrule
\endhead
\rlab{1. Product design} & Someone designs the analytics tool & The product may be
flawed\\
\rowcolor{tablealt}
\rlab{2. Tool setup} & The tool is installed and preferences are set & The setup
may be wrong --- mis-tagged pages, duplicate tracking, wrong filters\\
\rlab{3. Data output} & The tool generates reports & Sampling, bot traffic,
consent gaps\\
\rowcolor{tablealt}
\rlab{4. Human interpretation} & Someone reads the data & The interpretation may
be wrong --- correlation read as causation, averages hiding segments\\
\rlab{5. Business action} & Reports are used to make choices & The final decision
may be wrong\\
\end{mmlongtable}

\begin{keybox}[The Analyst's Real Job]
A number is not a fact and a report is not an insight. The analyst's value lies in
steps 2, 4 and 5 --- instrumenting correctly, interpreting honestly, and
translating the interpretation into a decision the business can act on.

\smallskip
An organisation with excellent tools and poor interpretation makes confident
wrong decisions \emph{faster}.
\end{keybox}

\section{How the Analyst Influences the Business}

\begin{enumerate}
\item \sfb{Sets the measurement agenda.} By defining which KPIs are reported, the
      analyst determines what the organisation pays attention to.
\item \sfb{Arbitrates between opinions.} Analytics ends debates that seniority
      would otherwise decide.
\item \sfb{Reallocates money.} Channel and campaign analysis directly redirects
      marketing budget.
\item \sfb{Prioritises the roadmap.} Funnel analysis determines which product and
      site fixes are built first.
\item \sfb{Quantifies risk and opportunity.} Forecasts and cohort models inform
      pricing, inventory and hiring.
\item \sfb{Protects against self-deception.} The analyst is the person who reports
      that the favourite campaign did not work.
\item \sfb{Translates.} The analyst converts technical measurement into business
      language that executives can act on --- and this, more than technical skill,
      determines their influence.
\end{enumerate}

\section{The Analyst as a Change Agent}

\begin{keybox}[The Syllabus Names This Explicitly]
Reporting numbers is not change; changing what the organisation \emph{does} is.
\end{keybox}

\begin{mmlongtable}{From reporting function to change agent}{tab:changeagent}%
{@{}P{0.45\textwidth}P{0.48\textwidth}@{}}
\rowcolor{ocre}\thd{From (reporting function)} & \thd{To (change agent)}\\
\midrule
\endfirsthead
\rowcolor{ocre}\thd{From (reporting function)} & \thd{To (change agent)}\\
\midrule
\endhead
Produces scheduled reports on request & Proactively identifies problems nobody
asked about\\
\rowcolor{tablealt}
Answers ``what happened?'' & Answers ``why did it happen, and what should we
do?''\\
Describes performance & Recommends a specific action with an expected outcome\\
\rowcolor{tablealt}
Owned by one team & Embedded across marketing, product, sales and service\\
Measures after the fact & Designs experiments before decisions are made\\
\rowcolor{tablealt}
Presents data & Builds the business case and drives adoption of the decision\\
\end{mmlongtable}

\subsection{What being a change agent requires}

\begin{itemize}
\item \sfb{Business literacy} --- understanding the commercial model well enough
      to know which questions matter.
\item \sfb{Credibility} --- data quality and honest reporting, including
      inconvenient findings.
\item \sfb{Storytelling} --- a finding that cannot be communicated will not change
      anything.
\item \sfb{Stakeholder relationships} --- influence flows through trust, not
      through dashboards.
\item \sfb{A bias towards experimentation} --- proposing tests rather than
      asserting opinions.
\item \sfb{Follow-through} --- measuring whether the implemented change actually
      worked, and saying so.
\end{itemize}

\subsection{Barriers a change agent must overcome}

\begin{warnbox}[The Six Barriers]
\begin{itemize}
\item \sfb{Data distrust} --- ``the numbers must be wrong'' is the standard first
      response to unwelcome findings.
\item \sfb{HiPPO decision-making} --- the \keytermas{highest paid person's
      opinion}{HiPPO} overriding evidence.
\item \sfb{Vanity metrics} --- organisations preferring numbers that always go up.
\item \sfb{Siloed ownership} of data across departments.
\item \sfb{Privacy and consent constraints} limiting what can be measured.
\item \sfb{Analysis paralysis} --- endless measurement substituting for decision.
\end{itemize}
\end{warnbox}

\section{Designing a Holistic Marketing Dashboard}

\begin{keybox}[High Yield]
``Design a basic digital dashboard that reflects holistic marketing
performance.''
\pyq{SEE, Jun/Jul 2025, Q10a --- 8 M}
A good dashboard answers three questions in order --- \emph{are we growing?},
\emph{which channels are working?}, and \emph{where are we losing people?} --- and
every metric on it must be tied to a decision someone can make.
\end{keybox}

\begin{mmlongtable}{A holistic marketing dashboard}{tab:dashboard}%
{@{}P{0.20\textwidth}P{0.44\textwidth}P{0.29\textwidth}@{}}
\rowcolor{ocre}\thd{Dashboard section} & \thd{Metrics} & \thd{Decision it supports}\\
\midrule
\endfirsthead
\rowcolor{ocre}\thd{Dashboard section} & \thd{Metrics} & \thd{Decision it supports}\\
\midrule
\endhead
\rlab{1. Executive summary} & Revenue, conversions, conversion rate, customer
acquisition cost, return on advertising spend, lifetime value against acquisition
cost --- each against target and prior period & Is the business on plan?\\
\rowcolor{tablealt}
\rlab{2. Acquisition} & Sessions and conversions by channel; spend, cost per click
and cost per acquisition per channel; new against returning & Where should the
next rupee of budget go?\\
\rlab{3. Engagement} & Bounce rate, session duration, pages per session, scroll
depth, best and worst landing pages & Which content and pages need work?\\
\rowcolor{tablealt}
\rlab{4. Conversion funnel} & Step-by-step drop-off from landing to purchase,
split by device & Which single step do we fix first?\\
\rlab{5. Retention and loyalty} & Repeat purchase rate, cohort retention curves,
churn, net promoter score, review sentiment & Are we keeping the customers we
buy?\\
\rowcolor{tablealt}
\rlab{6. Channel deep-dives} & SEO --- impressions, rankings, organic sessions;
e-mail --- open rate, click-through rate, revenue per message; social --- reach,
engagement rate, share of voice; mobile --- installs, active users, retention &
Which channel-level tactic needs attention?\\
\rlab{7. Technical health} & Core Web Vitals, error rate, uptime, tracking
integrity & Is anything silently broken?\\
\end{mmlongtable}

\begin{keybox}[Dashboard Design Principles]
One screen, top-down importance. Show \textbf{trend}, not just the current value
--- a number without a comparison is meaningless. Always include the
\textbf{target} alongside the actual. \textbf{Segment} by device and channel,
because averages hide failures. Use \textbf{colour only to signal status}, never
for decoration. And apply the same rule as web design: \textbf{if a metric drives
no decision, delete it.}
\end{keybox}

\section{Applying Digital Analytics --- Two Worked Scenarios}

\begin{exambox}[Improving Customer Retention for a Food Delivery App]
\pyq{CIE-III, Jun 2026 --- 10 M}
\textbf{Official scheme answer.} The app can analyse customer ordering patterns,
feedback and abandoned carts. Personalised offers, loyalty programmes and targeted
notifications can improve retention. Analytics helps identify customer needs and
improve services effectively.

\smallskip
\textbf{A fuller answer, in five steps.}
\begin{enumerate}
\item \sfb{Measure} --- build cohort retention curves by acquisition source, and
      calculate churn rate, order frequency and average order value per segment.
\item \sfb{Segment} --- identify the behaviours that distinguish retained from
      churned users, such as ordering within the first seven days or using a saved
      address.
\item \sfb{Diagnose} --- analyse cart abandonment reasons, delivery-time
      complaints and support tickets to locate the experience failure.
\item \sfb{Act} --- trigger personalised offers based on past cuisine preference;
      run a win-back push after 14 inactive days; introduce a subscription pass for
      free delivery; add a loyalty streak to build habit; and fix the delivery-time
      promise that the data shows is most often broken.
\item \sfb{Measure again} --- track Day-30 retention, repeat order rate and
      lifetime value against a control group, to prove the change worked.
\end{enumerate}
\end{exambox}

\begin{exambox}[Maintaining Competitive Advantage in a Global Market]
\pyq{CIE-III, Jun 2026 --- 10 M}
\textbf{Official scheme answer.} Digital analytics helps understand customer
preferences, competitor trends and campaign effectiveness. Website optimization
improves user experience and conversions. Together, they support better
decision-making, customer engagement and global competitiveness.

\smallskip
\textbf{A fuller answer.} Analytics provides market-by-market insight into what
customers actually want, which channels perform in each region, and where
competitors are gaining share. Website optimization ensures that the traffic won
in each market \emph{converts} --- through localisation of language, currency,
payment methods and page speed on regional networks. The combination creates a
continuous feedback loop in which the business learns faster than its competitors.

\smallskip
\textbf{Speed of learning, not size, is the durable advantage in a global digital
market.}
\end{exambox}

\begin{summarybox}
Data requires interpretation, and the chain from product design through tool
setup, data output and human interpretation to business action can break at every
step; the analyst's value is concentrated in instrumenting, interpreting and
translating. The analyst influences the business by setting the measurement
agenda, arbitrating between opinions, reallocating money, prioritising the
roadmap, quantifying risk, protecting against self-deception and translating
measurement into business language. Becoming a change agent means moving from
scheduled reports to proactive problem-finding, from ``what happened'' to ``what
should we do'', and from presenting data to driving adoption --- which requires
business literacy, credibility, storytelling, relationships, a bias towards
experimentation and follow-through, against barriers of data distrust,
HiPPO decision-making, vanity metrics, silos, privacy constraints and analysis
paralysis. A holistic dashboard has seven sections, each tied to a decision, and
any metric that drives no decision should be deleted.
\end{summarybox}

%%%%%%%%%%%%%%%%%%%%%%%%%%%%%%%%%%%%%%%%%%%%%%%%%%%%%%%%%%%%%%%%%%%%%%%%%%%%%%
%  Unit V appendices
%%%%%%%%%%%%%%%%%%%%%%%%%%%%%%%%%%%%%%%%%%%%%%%%%%%%%%%%%%%%%%%%%%%%%%%%%%%%%%

\startappendices

\chapter{Syllabus-to-Chapter Concordance}

This appendix exists so that coverage can be \emph{audited} rather than assumed.
The left-hand column reproduces the Unit~V syllabus descriptor phrase by phrase,
in the order the official document prints it. The right-hand column gives the
chapter and section that carries it. Every phrase resolves to a location, and no
numbered chapter of this volume contains material that is not traceable to a
phrase in this table.

\begin{mmlongtable}{Concordance for Unit V}{tab:concordance5}%
{@{}P{0.36\textwidth}P{0.57\textwidth}@{}}
\rowcolor{ocre}\thd{Syllabus phrase} & \thd{Where it is covered}\\
\midrule
\endfirsthead
\rowcolor{ocre}\thd{Syllabus phrase} & \thd{Where it is covered}\\
\midrule
\endhead
\multicolumn{2}{@{}l}{\bfseries\color{ocredark} First half --- Conversion Optimization}\\
\midrule
\rlab{What is AIDAS and its role} & Chapter~1 entire. What AIDAS stands for
\S\,1.1; the five stages \S\,1.2; why Satisfaction is critical \S\,1.3; the role
of AIDAS in CRO \S\,1.4; structuring a landing page with AIDAS \S\,1.5.\\
\rowcolor{tablealt}
\rlab{website optimization} & Chapter~2 entire. Conversion, conversion rate and
CRO defined \S\,2.1; why CRO matters \S\,2.2; the conversion funnel and its
drop-offs \S\,2.3; KPIs for website optimization \S\,2.4; the optimisation
workflow \S\,2.5.\\
\rlab{what visitors want to see on the website} & Chapter~3 entire. The ten
visitor questions \S\,3.1; the psychology beneath them \S\,3.2.\\
\rowcolor{tablealt}
\rlab{how to optimize key element and increase the effect of landing on a
particular page} & Chapter~4 entire. Headline \S\,4.1; call to action \S\,4.2;
visuals \S\,4.3; forms \S\,4.4; trust signals and social proof \S\,4.5; page
speed and mobile \S\,4.6.\\
\midrule
\multicolumn{2}{@{}l}{\bfseries\color{ocredark} Second half --- Digital Analytics}\\
\midrule
\rlab{Evolution of Digital Analytics} & Chapter~5 entire. Digital analytics
defined and distinguished from web analytics \S\,5.1; influence on business
decisions \S\,5.2; the six eras \S\,5.3.\\
\rowcolor{tablealt}
\rlab{information about end-to-end customer experience} & Chapter~6 entire. What
it means \S\,6.1; why the siloed view fails \S\,6.2; the seven requirements
\S\,6.3; business value \S\,6.4; refining segmentation \S\,6.5.\\
\rlab{analyst's influence on business} & Chapter~7, \S\,7.1 and \S\,7.2 --- the
interpretation chain and the seven modes of influence. The dashboard of \S\,7.5
and the worked scenarios of \S\,7.6 are that influence in practice.\\
\rowcolor{tablealt}
\rlab{role as a change agent} & Chapter~7, \S\,7.3 --- the shift from reporting
function to change agent, what it requires, and the six barriers.\\
\end{mmlongtable}

\begin{keybox}[Two Notes on Coverage]
\textbf{Deliberate overlap with Volume~4.} ``What visitors want to see'' is named
by both the Unit~IV descriptor (under user-centred design) and the Unit~V
descriptor. Volume~4 treats it as a design brief; Chapter~3 here treats it as a
conversion diagnosis and links it explicitly to the AIDAS\ix{AIDAS} stages.

\smallskip
\textbf{No off-syllabus appendix.} As in Volumes~3 and~4, every topic in the
source material for Unit~V maps to a descriptor phrase.
\end{keybox}

\chapter{Previous-Year Question Bank --- Unit V}

Every question in this appendix is reproduced from an actual RVCE Marketing
Management Continuous Internal Evaluation or Semester End Examination paper for
this course. Answers and pointers follow the official schemes of valuation where
those were released.

\section*{Part A --- Two-Mark Questions}
\addcontentsline{toc}{section}{Part A --- Two-Mark Questions}

\begin{enumerate}[leftmargin=2.2em,itemsep=1.1ex]

\item \textbf{What does the AIDAS model stand for in conversion optimization?}
      \pyqr{Improvement CIE, 4 Jun 2025 --- 2 M}\\
      Attention, Interest, Desire, Action, Satisfaction --- the sequence of
      psychological states a prospect passes through from first exposure to
      post-purchase satisfaction. \hfill\emph{See} \S\,1.1.

\item \textbf{Why is the ``Satisfaction'' component important in the AIDAS
      model?} \pyqr{SEE, Jun/Jul 2025 --- 2 M}\\
      Because the sale is not the end of the funnel. Satisfaction determines
      repurchase, retention and advocacy\ix{advocacy} --- and selling to an existing customer
      succeeds 60--70\% of the time against 5--20\% for a new prospect, with
      repeat customers spending up to 67\% more. A dissatisfied customer also
      broadcasts the failure publicly, damaging future acquisition.
      \hfill\emph{See} \S\,1.3.

\item \textbf{What is the role of AIDAS framework in the conversion
      optimization?} \pyqr{SEE, Oct/Nov 2023 --- 2 M}\\
      It acts as a diagnostic and structuring tool --- locating which
      psychological stage is failing, assigning each page section a defined job,
      sequencing the page accordingly, and generating testable optimisation
      hypotheses. \hfill\emph{See} \S\,1.4.

\item \textbf{What is digital analytics?} \pyqr{CIE-III, Jun 2026 --- 2 M}\\
      The collection, measurement, analysis and interpretation of digital data to
      improve business performance and customer experience.
      \hfill\emph{See} \S\,5.1.

\item \textbf{How does digital analytics influence business decisions?}
      \pyqr{CIE-III, Jun 2026 --- 2 M}\\
      It provides customer insights and performance data that help businesses
      improve marketing strategies and customer engagement --- determining budget
      allocation, product development, segment targeting, experience priorities
      and forecasts. \hfill\emph{See} \S\,5.2.

\item \textbf{What is meant by the ``end-to-end customer experience'' in digital
      analytics?} \pyqr{SEE, Jun/Jul 2025 --- 2 M}\\
      Measuring and understanding the complete customer journey\ix{customer journey} across every
      touchpoint and stage --- from first impression\ix{impression} through research, purchase,
      delivery and support to repeat purchase and advocacy --- as one connected
      experience rather than as disconnected channel reports.
      \hfill\emph{See} \S\,6.1.

\item \textbf{State two key elements that visitors commonly expect to see on a
      well-optimized landing page.} \pyqr{Improvement CIE, 4 Jun 2025 --- 2 M}\\
      A clear benefit-led headline that matches the advertisement or link that
      brought them there, and a single prominent call-to-action button. Also
      acceptable: trust signals such as testimonials or ratings; transparent
      pricing. \hfill\emph{See} \S\,3.1 and \S\,4.1.

\item \textbf{What is Website Optimization?} \pyqr{CIE-III, Jun 2026 --- 2 M}\\
      The process of improving website performance, speed, usability\ix{usability} and search
      visibility in order to increase traffic and conversions.
      \hfill\emph{See} \S\,2.1, and Volume~4 \S\,1.4.

\item \textbf{What is the impact of user behavior analysis on website
      optimization?} \pyqr{SEE, Oct/Nov 2023 --- 2 M}\\
      It reveals what users actually do rather than what is assumed ---
      identifying friction, ignored content, invisible calls to action, rage
      clicks and true scroll depth\ix{scroll depth} --- and produces the hypotheses that A/B testing
      then validates. \hfill\emph{See} \S\,2.3, and Volume~4 \S\,5.4.

\end{enumerate}

\section*{Part B --- Eight- and Ten-Mark Questions}
\addcontentsline{toc}{section}{Part B --- Eight- and Ten-Mark Questions}

\begin{enumerate}[leftmargin=2.2em,itemsep=1.4ex]

\item \textbf{Explain the AIDAS model and its role in website conversion
      optimization. How can key elements of a website be optimized to increase
      landing page effectiveness?} \pyqr{Improvement CIE, 4 Jun 2025 --- 10 M}\\
      \emph{Part one} --- define AIDAS and explain all five stages using the
      stage table, then give the five points on its role in CRO.
      \emph{Part two} --- take the key elements in turn (headline, call to action\ix{call to action},
      visuals, forms, trust signals, speed and mobile) and give the specific
      optimisation for each. Close with the AIDAS-structured landing page layout.
      \hfill\emph{See} Chapter~1 and Chapter~4.

\item \textbf{Explain the AIDAS model in detail. Discuss how each stage plays a
      critical role in conversion optimization, using examples to illustrate how
      businesses can enhance customer journeys through this model.}
      \pyqr{Improvement CIE, 28 Aug 2024 --- 10 M}\\
      As above, but weight the answer towards examples per stage.
      \emph{Attention} --- a streaming service's landing page whose entire first
      screen is one headline, one CTA and nothing else. \emph{Interest} --- a
      product page opening with three benefit bullets rather than technical
      specifications. \emph{Desire} --- a marketplace placing star ratings and
      review counts directly above the buy button. \emph{Action} --- a one-tap
      checkout with saved payment details and guest checkout enabled.
      \emph{Satisfaction} --- an order-confirmation and proactive
      delivery-tracking sequence followed by a review request, converting the
      buyer into an advocate. \hfill\emph{See} \S\,1.2 and \S\,1.5.

\item \textbf{Analyze how a company can use the AIDAS model to structure content
      on its product landing page.} \pyqr{SEE, Jun/Jul 2025, Q10b --- 8 M}\\
      The section-by-section table mapping each part of the page to an AIDAS
      stage, closing with the structuring principle that a section with no
      assignable stage should be deleted. \hfill\emph{See} \S\,1.5.

\item \textbf{Discuss what website visitors typically look for when they land on a
      business webpage. Suggest effective design and content strategies to meet
      those expectations.} \pyqr{Improvement CIE, 4 Jun 2025 --- 10 M}\\
      The ten visitor questions with the paired delivery for each, grounded in
      Krug's first law and the scanning behaviour that follows from it.
      \hfill\emph{See} Chapter~3.

\item \textbf{A travel agency has created a landing page for a limited-time offer
      on holiday packages. Despite a high number of clicks on the ads leading to
      the landing page, the conversion rate is disappointingly low. Evaluate the
      possible reasons for the low conversion rate on the landing page. Suggest
      optimizations for key elements such as headlines, call-to-action buttons,
      and visuals.} \pyqr{Improvement CIE, 28 Aug 2024 --- 10 M}\\
      \emph{Diagnose with AIDAS.} \textbf{Attention held} --- the clicks prove the
      advertisement worked, so the failure is after arrival. \textbf{Interest
      failing} --- poor message match\ix{message match} between the advertisement promise and the
      page headline; the offer is below the fold. \textbf{Desire failing} --- no
      testimonials, ratings, certifications or cancellation policy; pricing is
      opaque; no urgency despite the offer being time-limited; generic stock
      imagery. \textbf{Action failing} --- competing calls to action or none
      dominant; a long enquiry form; site navigation letting visitors wander away;
      slow mobile load.

      \emph{Optimisations --- headline:} restate the advertisement promise with
      specifics, for example ``Kerala backwaters, 4 nights from \rupee 18,999 ---
      book by Sunday'', with a sub-headline listing inclusions.
      \emph{CTA:} one dominant high-contrast button repeated above and below the
      fold, active copy (``Reserve my package''), reassurance microcopy (``Free
      cancellation up to 7 days'').
      \emph{Visuals:} real destination and customer photographs instead of stock, a
      short itinerary video, and a countdown timer.
      \emph{Additional:} cut the form to three fields, place testimonials adjacent
      to the CTA, remove navigation from the landing page, and compress images for
      a sub-three-second mobile load.
      \emph{Validate:} A/B test the headline and CTA, and confirm the fix through
      conversion rate, scroll depth and session recordings.
      \hfill\emph{See} \S\,1.4 and Chapter~4.

\item \textbf{Trace the evolution of digital analytics and its growing importance
      in delivering a seamless end-to-end customer experience.}
      \pyqr{Improvement CIE, 4 Jun 2025 --- 10 M}\\
      Give the six eras --- server logs, page-tag web analytics, visitor and goal
      analytics, multi-channel and social analytics, user-centric event-based
      analytics, and unified predictive privacy-first analytics --- with the tools
      and the limitation of each. Then explain the direction of travel (counting
      files \ra{} predicting behaviour; silos \ra{} unified view), and connect it
      to the end-to-end customer experience material, including the
      decision-constellation point and the seven requirements for genuine
      end-to-end measurement. \hfill\emph{See} \S\,5.3 and Chapter~6.

\item \textbf{A company plans to enter a competitive global market. Explain how
      digital analytics and website optimization can help maintain competitive
      advantage.} \pyqr{CIE-III, Jun 2026 --- 10 M}\\
      The official scheme answer expanded with market-by-market insight,
      localisation of language, currency, payment methods and regional page speed,
      and the argument that speed of learning is the durable advantage.
      \hfill\emph{See} \S\,7.6.

\item \textbf{A food delivery app wants to improve customer retention\ix{customer retention}. Apply
      digital analytics concepts to solve the problem.}
      \pyqr{CIE-III, Jun 2026 --- 10 M}\\
      The official scheme answer expanded into the five-step structure: measure
      \ra{} segment \ra{} diagnose \ra{} act \ra{} re-measure against a control
      group. \hfill\emph{See} \S\,7.6.

\item \textbf{Describe the conversion funnel\ix{conversion funnel} in digital marketing. What are the key
      stages and how can business identify and address drop offs in the funnel?}
      \pyqr{SEE, Oct/Nov 2023, Q9a --- 8 M}\\
      The five-stage funnel table with metrics and drop-off causes, the five
      identification methods, and the stage-by-stage remedies.
      \hfill\emph{See} \S\,2.3.

\item \textbf{Explain the concept of social proof and its impact on conversion
      rates. Provide examples of how business can use social proof effectively in
      marketing.} \pyqr{SEE, Oct/Nov 2023, Q9b --- 8 M}\\
      The definition, the impact argument grounded in the 90\%-against-48\% trust
      asymmetry, and the ten types of social proof with usage guidance --- plus the
      warning that fabricated proof destroys more trust than it borrows.
      \hfill\emph{See} \S\,4.5.

\item \textbf{Define Digital Analytics. Describe the key metrics and KPIs used in
      digital analytics. How can these metrics help business measure the success of
      their online campaigns?} \pyqr{SEE, Oct/Nov 2023, Q10a --- 10 M}\\
      Define digital analytics and distinguish it from web analytics. Present the
      KPIs in structured families --- \emph{acquisition} (sessions, users, traffic
      by channel, cost per click, cost per acquisition), \emph{engagement} (bounce\ix{bounce}
      rate, session duration, pages per session, scroll depth), \emph{conversion}
      (conversion rate, cart abandonment, form completion, revenue, average order
      value), \emph{retention} (repeat purchase rate, churn\ix{churn}, cohort retention, net
      promoter score, lifetime value), and \emph{efficiency} (return on advertising
      spend, lifetime value against acquisition cost).

      Then explain how they measure campaign success: they attribute results to
      specific campaigns and channels, allow comparison against targets and prior
      periods, expose which stage of the funnel a campaign fails at, and convert
      marketing activity into financial terms the business can act on.
      \hfill\emph{See} \S\,5.1 and \S\,2.4.

\item \textbf{What are the key performance indicators (KPIs) used to evaluate
      website optimization, and how do they contribute to increased conversions?}
      \pyqr{SEE, Jun/Jul 2025, Q9a --- 8 M}\\
      The twelve KPIs with the specific contribution each makes to conversion.
      \hfill\emph{See} \S\,2.4.

\item \textbf{Analyze how market segmentation can be refined using digital
      analytics data.} \pyqr{SEE, Jun/Jul 2025, Q9b --- 8 M}\\
      The seven refinements, from behavioural and value-based segmentation through
      to predictive segmentation and continuous automatic refresh.
      \hfill\emph{See} \S\,6.5.

\item \textbf{Design a basic digital dashboard that reflects holistic marketing\ix{holistic marketing}
      performance.} \pyqr{SEE, Jun/Jul 2025, Q10a --- 8 M}\\
      The seven-section dashboard with metrics and the decision each section
      supports, plus the design principles --- one screen, trend not just value,
      target alongside actual, segmented, colour for status only, and delete any
      metric that drives no decision. \hfill\emph{See} \S\,7.5.

\item \textbf{Demonstrate how tools like Google Analytics\ix{Google Analytics} and Hotjar can be used to
      develop insights into user behavior and improve marketing decisions.}
      \pyqr{SEE, Jun/Jul 2025, Q7a --- 8 M}\\
      Analytics supplies the \emph{what} --- channel performance, landing-page
      bounce rates, funnel drop-off, converting segments, cost per acquisition.
      Behavioural tools supply the \emph{why} --- heatmaps showing what attracts
      attention, scroll maps showing whether users reach the CTA, click maps
      exposing rage clicks on non-clickable elements, session recordings revealing
      hesitation, and on-page surveys capturing the abandonment reason in the
      user's own words.

      \emph{Combined workflow:} analytics shows the pricing page loses 60\% of
      visitors \ra{} recordings show users repeatedly clicking a static comparison
      table \ra{} the table is made interactive \ra{} an A/B test validates the
      change \ra{} analytics confirms improved funnel conversion. This is the
      analyst-as-change-agent loop applied in practice.
      \hfill\emph{See} \S\,7.3, and Volume~3 \S\,3.2.

\end{enumerate}

\section*{How Unit V Is Examined}
\addcontentsline{toc}{section}{How Unit V Is Examined}

\begin{keybox}[The Pattern Across Papers]
In the Semester End Examination, Unit~V is Questions~9 and~10 with internal
choice --- 16 marks in two eight-mark sub-divisions --- plus two or three Part~A
objective questions. In the Improvement CIEs it dominates alongside Unit~IV.

\smallskip
\textbf{AIDAS is effectively guaranteed.} It has appeared in every Unit~V paper
available, as both a two-mark definition and a ten-mark explanation. The other
high-probability topics are \textbf{landing page optimisation} for a
low-converting page, the \textbf{evolution of digital analytics}, the
\textbf{conversion funnel} and drop-off analysis, \textbf{social proof}, and
\textbf{KPIs and dashboards}.

\smallskip
As in Unit~IV, most Part-B questions are scenarios. Answer them with the AIDAS
diagnostic structure, or with the measure \ra{} diagnose \ra{} hypothesise \ra{}
test \ra{} implement \ra{} re-measure loop --- not with general commentary.
\end{keybox}

\chapter{Glossary of Key Terms}

\begin{mmlongtable}{Key terms for rapid revision}{tab:glossary5}%
{@{}P{0.26\textwidth}P{0.67\textwidth}@{}}
\rowcolor{ocre}\thd{Term} & \thd{One-line meaning}\\
\midrule
\endfirsthead
\rowcolor{ocre}\thd{Term} & \thd{One-line meaning}\\
\midrule
\endhead
\rlab{A/B test} & Splitting traffic between two versions to find the statistically
better one\\
\rowcolor{tablealt}
\rlab{Above the fold} & The area visible without scrolling\\
\rlab{AIDA} & The classical four-stage model: Attention, Interest, Desire,
Action\\
\rowcolor{tablealt}
\rlab{AIDAS} & Attention, Interest, Desire, Action, Satisfaction\\
\rlab{Cart abandonment rate} & Percentage who begin checkout but do not complete
it\\
\rowcolor{tablealt}
\rlab{CDP} & Customer data platform --- unifies customer data into one profile\\
\rlab{Change agent} & An analyst who drives decisions and organisational change,
not just reports\\
\rowcolor{tablealt}
\rlab{Churn rate} & Percentage of customers lost in a period\\
\rlab{Cohort analysis} & Tracking a group who share a start period over time\\
\rowcolor{tablealt}
\rlab{Conversion} & A defined valuable action completed by a visitor\\
\rlab{Conversion funnel} & The staged path from awareness to purchase and
retention\\
\rowcolor{tablealt}
\rlab{Conversion rate} & Conversions divided by total visitors, times 100\\
\rlab{CRO} & Systematically increasing the percentage of visitors who convert,
using data and experiments\\
\rowcolor{tablealt}
\rlab{Decision constellation} & The non-linear network of simultaneous
micro-decisions replacing the classic funnel\\
\rlab{Digital analytics} & Collection, measurement, analysis and interpretation of
digital data to improve business performance and customer experience\\
\rowcolor{tablealt}
\rlab{Drop-off} & The proportion of users lost between two funnel steps\\
\rlab{End-to-end customer experience} & The complete journey across all
touchpoints, measured as one connected experience\\
\rowcolor{tablealt}
\rlab{HiPPO} & ``Highest paid person's opinion'' --- the bias analytics exists to
counter\\
\rlab{Landing page} & A single-purpose page built for one conversion\\
\rowcolor{tablealt}
\rlab{Last-click attribution} & Crediting the final touchpoint only --- the
classic distortion\\
\rlab{LTV : CAC} & Lifetime value against customer acquisition cost --- the core
efficiency ratio\\
\rowcolor{tablealt}
\rlab{Message match} & Consistency between the advertisement or link and the
landing page it leads to\\
\rlab{Multi-touch attribution} & Distributing conversion credit across all
contributing touchpoints\\
\rowcolor{tablealt}
\rlab{NPS} & Net promoter score --- likelihood of recommendation\\
\rlab{RFM analysis} & Segmenting by recency, frequency and monetary value\\
\rowcolor{tablealt}
\rlab{Risk reversal} & A guarantee or free-return policy that removes the buyer's
perceived risk\\
\rlab{Social proof} & People assuming others' actions reflect correct behaviour\\
\rowcolor{tablealt}
\rlab{Vanity metric} & A number that always rises and drives no decision\\
\end{mmlongtable}

\chapter{Framework Quick-Reference Sheet}

Everything in this volume that is a numbered list, in one place, for the night
before the examination.

\begin{mmlongtable}{Framework quick reference}{tab:quickref5}%
{@{}P{0.28\textwidth}P{0.65\textwidth}@{}}
\rowcolor{ocre}\thd{Framework} & \thd{The sequence}\\
\midrule
\endfirsthead
\rowcolor{ocre}\thd{Framework} & \thd{The sequence}\\
\midrule
\endhead
\rlab{AIDAS} & Attention \ra{} Interest \ra{} Desire \ra{} Action \ra{}
Satisfaction\\
\rowcolor{tablealt}
\rlab{AIDAS states} & ``caught my eye'' \ra{} ``might be relevant'' \ra{} ``I want
this'' \ra{} ``I am buying'' \ra{} ``I am glad I did''\\
\rlab{AIDAS diagnostic} & Immediate bounce = Attention; bounce after scroll =
Interest; add-to-cart without checkout = Action; refunds and no repeat =
Satisfaction\\
\rowcolor{tablealt}
\rlab{Satisfaction economics} & Existing customer 60--70\% against new prospect
5--20\%; repeat customers spend up to 67\% more\\
\rlab{AIDAS role in CRO (five points)} & Locates the failure, maps content to
state, sequences the page, generates hypotheses, prevents over-optimising one
stage\\
\rowcolor{tablealt}
\rlab{AIDAS landing page} & Hero \ra{} value proposition \ra{} proof \ra{} offer
and urgency \ra{} conversion block \ra{} reassurance footer\\
\rlab{Conversion rate} & Conversions divided by total visitors, times 100\\
\rowcolor{tablealt}
\rlab{Why CRO matters} & Compounds on future traffic, cheaper than traffic,
improves every channel's ROI, lowers CAC, replaces opinion with evidence\\
\rlab{Conversion funnel} & Awareness \ra{} interest \ra{} desire \ra{} action
\ra{} satisfaction\\
\rowcolor{tablealt}
\rlab{Identifying drop-offs} & Build the funnel \ra{} segment it \ra{} diagnose
qualitatively \ra{} ask the user \ra{} benchmark\\
\rlab{Twelve CRO KPIs} & Conversion rate, bounce, duration and pages, load time,
exit rate, cart abandonment, form completion, CTA click-through, scroll depth,
mobile gap, return and repeat rate, CPA and ROAS\\
\rowcolor{tablealt}
\rlab{CRO workflow} & Measure \ra{} diagnose \ra{} hypothesise \ra{} prioritise
\ra{} test \ra{} implement \ra{} re-measure\\
\rlab{Ten visitor questions} & Right place, what's in it for me, trust, cost, next
step, findability, speed, mobile, payment safety, human contact\\
\rowcolor{tablealt}
\rlab{Headline rules} & Message match, specific and quantified, benefit-led,
short and dominant, plus a sub-headline\\
\rlab{CTA rules} & One dominant, high contrast, active first-person copy, above
the fold and repeated, 44 pixel tap target, reassurance microcopy\\
\rowcolor{tablealt}
\rlab{Form rules} & Drop unused fields, guest checkout, smart defaults, inline
validation, progress indicator, correct mobile keyboard\\
\rlab{Ten types of social proof} & Testimonials, ratings and reviews, user counts,
case studies, client and media logos, expert endorsement, influencer proof,
user-generated content, real-time activity, wisdom of the crowd\\
\rowcolor{tablealt}
\rlab{Six eras of digital analytics} & Server logs \ra{} page-tag web analytics
\ra{} visitor and goal analytics \ra{} multi-channel and social \ra{} user-centric
event-based \ra{} unified, predictive, privacy-first\\
\rlab{Direction of travel} & Counting files \ra{} visits \ra{} conversions \ra{}
understanding people \ra{} predicting behaviour; silos \ra{} unified view\\
\rowcolor{tablealt}
\rlab{End-to-end requirements (seven)} & Unified identity, consistent event
taxonomy, journey analysis, multi-touch attribution, voice-of-customer,
post-purchase data, consent governance\\
\rlab{Segmentation refinements (seven)} & Behavioural, value-based (RFM),
journey-stage, intent-based, channel-affinity, predictive, continuously
refreshed\\
\rowcolor{tablealt}
\rlab{Interpretation chain} & Product design \ra{} tool setup \ra{} data output
\ra{} human interpretation \ra{} business action\\
\rlab{Analyst's influence (seven)} & Sets the agenda, arbitrates, reallocates
money, prioritises the roadmap, quantifies risk, prevents self-deception,
translates\\
\rowcolor{tablealt}
\rlab{Change agent requires} & Business literacy, credibility, storytelling,
relationships, bias to experimentation, follow-through\\
\rlab{Six barriers} & Data distrust, HiPPO, vanity metrics, silos, privacy
constraints, analysis paralysis\\
\rowcolor{tablealt}
\rlab{Dashboard sections (seven)} & Executive summary, acquisition, engagement,
conversion funnel, retention and loyalty, channel deep-dives, technical health\\
\rlab{Dashboard principles} & One screen, show trend, include target, segment,
colour for status only, delete metrics that drive no decision\\
\end{mmlongtable}


%%%%%%%%%%%%%%%%%%%%%%%%%%%%%%%%%%%%%%%%%%%%%%%%%%%%%%%%%%%%%%%%%%%%%%%%%%%%%%
\backmatter
\printindex

\end{document}
